% ============================================================
% Chapter 2 — Background  &  Chapter 3 — Methodology
% ------------------------------------------------------------
% Standalone chapter file for the thesis
% "Evaluating multilingual LLM performance with cross-lingual alignment"
%
% Usage: % ============================================================
% Chapter 2 — Background  &  Chapter 3 — Methodology
% ------------------------------------------------------------
% Standalone chapter file for the thesis
% "Evaluating multilingual LLM performance with cross-lingual alignment"
%
% Usage: % ============================================================
% Chapter 2 — Background  &  Chapter 3 — Methodology
% ------------------------------------------------------------
% Standalone chapter file for the thesis
% "Evaluating multilingual LLM performance with cross-lingual alignment"
%
% Usage: % ============================================================
% Chapter 2 — Background  &  Chapter 3 — Methodology
% ------------------------------------------------------------
% Standalone chapter file for the thesis
% "Evaluating multilingual LLM performance with cross-lingual alignment"
%
% Usage: \input{chapters_2_3} from a report/book-class document
% (e.g. the tum-thesis-latex template). Requires: amsmath, amssymb,
% booktabs, natbib, and the accompanying references.bib.
% ============================================================


% ============================================================
% CHAPTER 2 — BACKGROUND
% ============================================================
\chapter{Background}
\label{ch:background}

This chapter introduces the concepts on which the thesis builds. We first
review how multilingual capabilities arise in large language models (LLMs)
and where they break down (Section~\ref{sec:multilingual-lms}). We then
discuss the special role of English as a pivot language, both as an explicit
engineering device in classical multilingual NLP and as an implicit property
of the internal representations of English-centric LLMs
(Section~\ref{sec:english-pivot}). Section~\ref{sec:multilingual-eval}
surveys how multilingual ability is evaluated today and motivates
alignment-based, reference-free metrics such as MEXA. Finally,
Section~\ref{sec:parallel-corpora} describes the parallel corpora that make
such evaluations possible and the resource disparities that constrain them.

% ------------------------------------------------------------
\section{Multilingual Language Models}
\label{sec:multilingual-lms}

% ..............................................
\subsection{Architectures and Pre-training}
\label{subsec:architectures}

Virtually all modern language models are based on the Transformer
architecture \citep{vaswani2017attention}, a stack of identical layers that
combine multi-head self-attention with position-wise feed-forward networks.
Each layer transforms a sequence of token representations into a new
sequence of contextualized representations; the output of layer $l$ for
token $t$, the \emph{hidden state} $h_t^{(l)}$, is the central object of
study in this thesis, because it is from these layer-wise hidden states
that we extract sentence embeddings.

Two architectural families dominate multilingual NLP. \emph{Encoder-only}
models such as mBERT \citep{devlin2019bert} and XLM-R
\citep{conneau2020xlmr} use bidirectional attention and are pre-trained
with masked-language-modeling objectives on deliberately multilingual
corpora (XLM-R covers 100 languages sampled from CommonCrawl). Their
representations are designed to serve downstream classification and
retrieval tasks, and dedicated sentence encoders such as LaBSE
\citep{feng2022labse}, LASER3 \citep{heffernan2022laser3}, and multilingual
E5 \citep{wang2024multilingual} additionally fine-tune on parallel data
with contrastive or translation-ranking objectives so that translations map
to nearby points in embedding space. \emph{Decoder-only} models such as
GPT-3 \citep{brown2020gpt3} and the Llama family \citep{dubey2024llama3}
use causal (left-to-right) attention and are pre-trained autoregressively
to predict the next token. They are the dominant paradigm for generative
LLMs, but their pre-training corpora are heavily skewed toward English:
for Llama~3, roughly 90\% of the pre-training tokens are English
\citep{dubey2024llama3}. Models of this kind are commonly called
\emph{English-centric}, even though they exhibit non-trivial abilities in
dozens of languages.

This asymmetry creates the central question of the thesis: an
English-centric decoder model is never explicitly trained to align
languages, yet it acquires multilingual competence as a by-product of
incidental multilingual data in its corpus. Understanding \emph{where} and
\emph{how well} this competence materializes inside the network --- rather
than only measuring it behaviorally on downstream tasks --- is what
alignment-based evaluation aims to do.

% ..............................................
\subsection{The Tokenization Bottleneck}
\label{subsec:tokenization}

Before any text reaches the Transformer layers, it is segmented into
subword units by a tokenizer, typically trained with byte-pair encoding
\citep{sennrich2016bpe} or a unigram language model as implemented in
SentencePiece \citep{kudo2018sentencepiece}. The tokenizer's vocabulary is
learned on (a sample of) the pre-training corpus, and therefore inherits
its language distribution: frequent English words receive dedicated tokens,
while words in under-represented languages are shattered into many short
fragments, in the worst case into individual bytes.

The severity of this effect is commonly quantified by \emph{fertility},
the average number of subword tokens produced per word (or per character
for scripts without whitespace word boundaries). High fertility has three
practical consequences. First, it inflates sequence lengths, making
inference disproportionately expensive for low-resource languages. Second,
it dilutes the semantic content carried by each token: a model that sees
Amharic or Tibetan only as byte sequences must reconstruct lexical meaning
from much noisier units than it does for English. Third, and most relevant
to this thesis, it affects the quality of sentence embeddings built from
token hidden states, since averaging over many weakly informative fragments
yields blurrier sentence representations. Non-Latin scripts (Ge'ez,
Myanmar, Khmer, Cherokee) and low-resource languages regardless of script
are the most affected. Throughout our experiments we therefore track
tokenizer fertility on both evaluation corpora alongside the alignment
scores, which allows us to distinguish genuinely weak cross-lingual
representations from representations that are handicapped already at the
tokenization stage.

% ..............................................
\subsection{Cross-lingual Transfer}
\label{subsec:transfer}

The practical value of multilingual models rests on \emph{cross-lingual
transfer}: capabilities acquired predominantly in one language (in
practice, English) become available in other languages without
language-specific supervision. The standard explanation is that the model
learns a partially \emph{shared semantic space} in which translation
equivalents --- words or sentences with the same meaning --- receive
similar representations, so that task knowledge attached to a region of
that space is accessible from any language that maps into it
\citep{hammerl2024survey}.

Empirical evidence supports this picture with an important refinement:
the degree of sharing is \emph{layer-dependent}. Input embeddings and the
first layers remain largely language-specific, dominated by lexical and
orthographic features; representations become increasingly
language-neutral through the middle of the network; and the final layers
re-specialize toward the output vocabulary of the target language.
Interpretability studies on Llama-family models show precisely this
trajectory \citep{wendler2024llamas,zhao2024multilingualism}, and it
motivates the layer-wise design of MEXA: if cross-lingual alignment exists
anywhere in an English-centric model, it should be found in the
intermediate layers, and a single-layer (e.g., last-layer) probe would
systematically underestimate it.

\citet{hammerl2024survey} organize the methods that \emph{create} or
\emph{improve} such alignment into families --- joint multilingual
pre-training, alignment objectives on parallel data, contrastive
sentence-level training, and post-hoc mapping between monolingual spaces.
This thesis is concerned with the complementary problem: not producing
alignment, but \emph{measuring} the alignment that a given model already
has, layer by layer and language by language.

% ------------------------------------------------------------
\section{English as a Pivot Language}
\label{sec:english-pivot}

% ..............................................
\subsection{Pivot Translation and Alignment}
\label{subsec:pivot-translation}

Using English as a bridge between languages long predates neural LLMs. In
statistical machine translation, direct parallel data between two
non-English languages is often scarce, so translation is routed through a
\emph{pivot} language: source $\rightarrow$ English $\rightarrow$ target
\citep{wu2007pivot}. Because English participates in the largest number of
parallel corpora, it is the pivot of choice, and the quality of any
source--target system becomes bounded by the quality of the two
English-anchored legs. Neural machine translation reproduced the pattern
in a softer form: multilingual NMT systems trained on English-centric data
develop an ``interlingua-like'' internal representation that enables
zero-shot translation between language pairs never seen together in
training \citep{johnson2017zeroshot}.

Modern decoder-only LLMs internalize this bridge. Using the logit-lens
technique, \citet{wendler2024llamas} show that when Llama-2 models process
non-English input, the intermediate layers pass through a phase in which
the representations are closest to \emph{English} tokens before the final
layers rotate them toward the target-language vocabulary --- the model
demonstrably ``works in English'' in its middle layers.
\citet{zhao2024multilingualism} corroborate this for reasoning tasks,
finding that multilingual inputs are converted into English-centric latent
representations in which the actual computation takes place. This
\emph{pivot hypothesis} is the theoretical foundation of MEXA: if
non-English understanding is routed through English-like intermediate
representations, then the degree to which a language's representations
align with English is a direct measurement of how much of the model's
(English-acquired) capability that language can access.

% ..............................................
\subsection{English Centricity and Bias}
\label{subsec:english-bias}

The pivot mechanism is a double-edged sword. On one side, it is exactly
what makes cross-lingual transfer from a 90\%-English corpus possible at
all. On the other, it builds a structural asymmetry into the model.
Languages that map cleanly onto English-like representations inherit the
model's full capability; languages that do not --- because of typological
distance, script, or sheer data scarcity --- are effectively locked out,
regardless of how much task knowledge the model possesses. Concepts and
distinctions that English does not encode (grammatical evidentiality,
honorific systems, culture-specific terminology) risk being flattened when
they are forced through an English-shaped bottleneck, and cultural framing
in generated text tends to default to Anglophone norms.

For evaluation, English centricity raises a methodological question that
this thesis addresses empirically: is English \emph{privileged} as a
measurement anchor, or merely \emph{convenient}? MEXA, as proposed, always
measures alignment against English. If the pivot hypothesis is the true
mechanism, replacing English with another pivot language should
systematically lower the measured alignment, in proportion to how far the
alternative pivot is from the model's internal hub. Our non-English-pivot
experiments (Section~\ref{sec:experimental-settings}) test exactly this by
recomputing alignment against French, German, Arabic, Chinese, and Basque
pivots.

% ------------------------------------------------------------
\section{Multilingual Evaluation}
\label{sec:multilingual-eval}

% ..............................................
\subsection{Standard Benchmarks}
\label{subsec:benchmarks}

The most direct way to assess multilingual capability is to run the model
on parallel downstream tasks. FLORES-200 \citep{nllbteam2022} provides
professionally translated Wikimedia sentences in 204 language--script
pairs and is the de-facto standard for evaluating translation and
cross-lingual retrieval. Belebele \citep{bandarkar2024belebele} builds a
machine reading comprehension benchmark on top of FLORES passages,
covering 122 language variants with fully parallel multiple-choice
questions --- crucially, differences in scores across languages can only
stem from the language itself, not from differing content. Knowledge- and
reasoning-oriented English benchmarks such as MMLU
\citep{hendrycks2021mmlu} and ARC have been machine-translated into
26--31 languages to form m-MMLU and m-ARC \citep{lai2023okapi}.

These benchmarks are indispensable, but they have structural limits: they
exist for at most a couple of hundred languages (Belebele: 122, m-MMLU:
34, m-ARC: 31), they are expensive to extend because each new language
needs high-quality translated task data, and machine-translated variants
inherit translation noise. For the several thousand remaining languages,
task-based evaluation is simply unavailable --- which is the gap that
alignment-based estimation aims to fill. In this thesis, downstream
benchmarks play the role of \emph{validation targets}: MEXA scores are
correlated against Belebele and m-ARC performance to test whether
alignment is a faithful proxy for task capability.

% ..............................................
\subsection{Reference-Based vs.\ Reference-Free Metrics}
\label{subsec:metrics}

Evaluation metrics in multilingual NLP fall into two broad families.
\emph{Reference-based} metrics compare model output against gold
references. Surface-overlap metrics such as BLEU
\citep{papineni2002bleu} are cheap and language-agnostic in principle, but
they reward n-gram overlap rather than meaning, penalize legitimate
paraphrases, and interact badly with morphologically rich languages and
non-whitespace scripts. Learned metrics such as COMET \citep{rei2020comet}
correlate far better with human judgments by scoring
source--hypothesis--reference triples with a fine-tuned multilingual
encoder, but they require human references and a trained scoring model,
and their reliability degrades precisely for the low-resource languages
where evaluation is most needed.

\emph{Reference-free} (quality-estimation or model-based) metrics remove
the dependency on gold outputs. Embedding-similarity approaches compare
representations of source and translation directly
\citep{reimers2020multilingual}, and model-internal approaches --- the
category MEXA belongs to --- do not look at generated text at all.
Instead, they probe the model's hidden states on parallel \emph{inputs}
and quantify structure that is known to correlate with capability. The
advantages are substantial: no task data, no references, no additional
trained judge, and applicability to any language for which a hundred
parallel sentences exist. The corresponding risk is validity --- an
internal signal is only useful insofar as it predicts external behavior
--- which is why correlation with downstream benchmarks
(Section~\ref{subsec:benchmarks}) is an essential part of the methodology
rather than an afterthought.

% ..............................................
\subsection{Challenges in Cross-lingual Alignment Evaluation}
\label{subsec:alignment-challenges}

Measuring alignment from raw embedding geometry is less straightforward
than it appears, for reasons rooted in the geometry of learned
representation spaces.

\paragraph{Anisotropy.} Transformer embedding spaces are strongly
anisotropic: representations occupy a narrow cone rather than spreading
over the full space, so the cosine similarity between two \emph{unrelated}
sentences is already high \citep{rajaee2022isotropy}. In multilingual
models the effect is compounded by outlier dimensions with extreme
variance that dominate cosine computations; removing a handful of such
dimensions can change similarity rankings entirely
\citep{hammerl2023anisotropy}. Absolute cosine values are therefore not
comparable across models, layers, or languages.

\paragraph{Hubness.} In high-dimensional spaces, some points become
``hubs'' that appear among the nearest neighbors of disproportionately
many queries. A naive nearest-neighbor criterion over cosine similarities
rewards such hubs and distorts retrieval-based measurements.

\paragraph{Language identity vs.\ semantics.} Embeddings encode both
\emph{what} a sentence says and \emph{which language} it is written in.
Two sentences in the same language are often closer to each other than a
sentence and its translation, simply because of the shared language
signal. A useful alignment metric must isolate the semantic component
from this language-identity component.

MEXA's design responds to all three problems at once: instead of using
cosine similarity as an absolute quantity, it uses it only
\emph{comparatively}, asking whether the correct translation outranks all
distractors within the same similarity matrix (a bidirectional
precision-at-1 criterion, formalized in
Section~\ref{sec:mexa-pipeline}). Rank-based comparisons within a fixed
matrix are invariant to the global scale and offset distortions that
anisotropy introduces, and the bidirectional requirement suppresses hub
sentences that would pass a one-directional test. As an extension, this
thesis additionally evaluates a \emph{mean-centered} variant in which each
language's mean embedding is subtracted before computing similarities,
directly removing the language-identity component discussed above.

% ------------------------------------------------------------
\section{Parallel Corpora}
\label{sec:parallel-corpora}

% ..............................................
\subsection{Sources and Typology}
\label{subsec:corpora-sources}

Alignment-based evaluation requires \emph{parallel} data: the same
sentences translated across many languages, so that semantic content is
held constant while language varies. Parallel corpora differ along
several axes that matter for measurement quality: translation quality
(professional vs.\ crawled/mined, e.g.\ OPUS \citep{tiedemann2012opus}),
domain (news, web, religious text), degree of parallelism (bitexts vs.\
$n$-way parallel corpora), and language coverage.

Two corpora sit at opposite ends of the coverage--quality trade-off and
are both used in this thesis. \textbf{FLORES-200} \citep{nllbteam2022}
contains professionally translated, quality-assured sentences from
Wikimedia sources in 204 language--script pairs; it is $n$-way parallel,
modern-domain, and high quality, but its coverage stops at roughly 200
languages. The \textbf{Parallel Bible Corpus} \citep{mayer2014bible}
exploits the fact that the Bible is the most translated text in existence,
with verse-level alignment across more than 1{,}400 language--script
pairs; the derived \emph{super-parallel} subset (sPBC) restricts to verses
available in all covered languages. Its domain is narrow and archaic and
translation styles vary, but it is the only resource that reaches deep
into the low-resource tail, which is exactly where reference-free
evaluation is most valuable.

% ..............................................
\subsection{Resource Disparity}
\label{subsec:resource-disparity}

The distribution of parallel (and monolingual) data across the world's
languages is extremely skewed. A few dozen high-resource languages account
for the overwhelming majority of available text; thousands of languages
have effectively no digital footprint beyond a Bible translation. This
disparity propagates through the entire pipeline studied in this thesis:
low-resource languages receive less pre-training data (weaker
representations), fragmented tokenization
(Section~\ref{subsec:tokenization}), and no downstream benchmarks
(Section~\ref{subsec:benchmarks}) --- so their capability can be neither
learned as well nor measured as easily.

For evaluation design, resource disparity has two concrete consequences.
First, any claim about ``multilingual'' capability that is validated only
on high-resource languages is systematically optimistic; extending
measurement to the sPBC's low-resource tail is therefore a substantive
robustness test, not an ornament. Second, corpus effects must be
disentangled from language effects: a language may score lower on the
Bible corpus than on FLORES because the model is genuinely weaker in it,
or because archaic religious register is far from the model's training
distribution. Using both corpora on the same models, as done here, makes
this confound visible and quantifiable.


% ============================================================
% CHAPTER 3 — METHODOLOGY
% ============================================================
\chapter{Methodology}
\label{ch:methodology}

This chapter specifies the experimental  of the thesis. We first
formalize the MEXA pipeline as implemented in our codebase
(Section~\ref{sec:mexa-pipeline}), then describe the evaluated models
(Section~\ref{sec:models}), the parallel corpora and their subsets
(Section~\ref{sec:datasets}), and finally the experimental settings,
analysis methods, and compute environment
(Section~\ref{sec:experimental-settings}). The methodology follows
\citet{kargaran2025mexa} for the reproduction experiments and extends it
along four axes: new decoder families, mixture-of-experts architectures,
encoder and embedding models, and non-English pivot languages.

% ------------------------------------------------------------
\section{The MEXA Pipeline}
\label{sec:mexa-pipeline}

MEXA (Multilingual Evaluation via Cross-Lingual Alignment) estimates the
multilingual coverage of a language model $m$ from parallel sentences
alone, in four stages: sentence-embedding extraction, layer-wise
similarity computation, alignment scoring, and aggregation.

\paragraph{Stage 1: Layer-wise sentence embeddings.}
Let $\mathcal{S} = \{s_1, \dots, s_n\}$ be a set of $n$ parallel sentences
available in a target language $L_1$ and a pivot language $L_2$ (English,
unless stated otherwise). Each sentence is passed through the model, and
hidden states $h_t^{(l)}$ are collected for every token $t$ at every layer
$l \in \{1, \dots, L\}$, including the embedding layer. Token states are
reduced to a single sentence embedding $e^{(l)}$ per layer with one of two
strategies:

\begin{itemize}
    \item \textbf{Position-weighted average.} Because causal attention
    means early tokens see little context while late tokens see all of it,
    tokens are averaged with weights proportional to their position:
    \begin{equation}
        e^{(l)} \;=\; \sum_{t=1}^{T} w_t\, h_t^{(l)},
        \qquad
        w_t \;=\; \frac{t}{\sum_{k=1}^{T} k},
        \label{eq:weighted-avg}
    \end{equation}
    where $T$ is the number of non-padding tokens. This is the primary
    setting in all reported results.
    \item \textbf{Last token.} The hidden state of the final non-padding
    token, $e^{(l)} = h_T^{(l)}$, which under causal attention is the only
    position that has attended to the entire sentence. This setting is
    used as a robustness check.
\end{itemize}

For encoder-only models (Section~\ref{subsec:encoder-models}), whose
attention is bidirectional, the same two reductions are applied to the
encoder hidden states, making the pipeline architecture-agnostic.
Embeddings are stored per language and per layer, so that scoring is
decoupled from (and far cheaper than) extraction.

\paragraph{Stage 2: Similarity matrix.}
For each layer $l$, a cosine similarity matrix
$C(L_1, L_2, m, l) \in \mathbb{R}^{n \times n}$ is computed between the
$n$ target-language embeddings and the $n$ pivot-language embeddings:
$c_{ij}(l)$ is the cosine similarity between target sentence $i$ and pivot
sentence $j$. The diagonal entries $c_{ii}(l)$ correspond to true parallel
pairs; all off-diagonal entries are non-parallel distractors.

\paragraph{Stage 3: Alignment score.}
As discussed in Section~\ref{subsec:alignment-challenges}, absolute cosine
values are unreliable in anisotropic LLM embedding spaces. MEXA therefore
scores alignment with a \emph{bidirectional retrieval} (precision-at-1)
criterion: a sentence pair counts as aligned only if its true similarity
strictly exceeds every distractor in both its row and its column,

\begin{equation}
    \mu C(L_1, L_2, m, l)
    \;=\;
    \frac{1}{n}
    \sum_{i=1}^{n}
    \mathbb{1}\!\left(
        c_{ii}(l) \;>\; \max_{j \neq i}\,\{\, c_{ij}(l),\; c_{ji}(l) \,\}
    \right).
    \label{eq:mexa}
\end{equation}

The score lies in $[0, 1]$; the expected value under random embeddings is
approximately $1/n^2$ per pair (both the row and the column condition must
hold by chance), so with $n = 100$ sentences a random baseline scores
essentially zero, which validates the use of small parallel samples.

\paragraph{Stage 4: Aggregation across layers.}
The layer-wise scores are summarized into two model-level statistics:
\begin{equation}
    \mu_{\text{Max}}(L_1) = \max_{l} \; \mu C(L_1, L_2, m, l),
    \qquad
    \mu_{\text{Mean}}(L_1) = \frac{1}{L} \sum_{l=1}^{L} \mu C(L_1, L_2, m, l).
\end{equation}
$\mu_{\text{Max}}$ captures the peak alignment a model achieves anywhere
in its depth --- typically in the middle-to-late layers, consistent with
the pivot hypothesis --- while $\mu_{\text{Mean}}$ rewards models that
sustain alignment across many layers. Model-level scores reported in
Chapter~4 average these quantities over all evaluated languages.

\paragraph{Extension: mean-centered scoring.}
To separate semantic alignment from the language-identity signal
(Section~\ref{subsec:alignment-challenges}), we additionally compute a
\emph{centered} variant of Equation~\eqref{eq:mexa}: before building the
similarity matrix at layer $l$, each language's embeddings are shifted by
their own mean,
\begin{equation}
    \tilde{e}_i^{(l)} \;=\; e_i^{(l)} - \frac{1}{n}\sum_{k=1}^{n} e_k^{(l)},
\end{equation}
applied independently to the target and pivot sides. Centering removes the
dominant direction shared by all sentences of a language and typically
sharpens retrieval; comparing centered and uncentered scores quantifies
how much measured alignment is carried by language-mean offsets rather
than sentence-level semantics.

\paragraph{Extension: non-English pivots.}
Equation~\eqref{eq:mexa} does not privilege English mathematically; the
pivot $L_2$ is a free parameter. We exploit this to test the pivot
hypothesis directly, recomputing all alignment scores with French, German,
Arabic, Chinese, and Basque as pivot languages --- spanning high-resource
Latin-script, non-Latin-script, and a language isolate --- in both the
standard and centered variants.

% ------------------------------------------------------------
\section{Models}
\label{sec:models}

We evaluate 25+ publicly available checkpoints, grouped into four
categories that correspond to the four experimental axes of the thesis.
All models are obtained from the Hugging Face Hub and evaluated without
any fine-tuning; base (non-instruction-tuned) variants are used where
available so that scores reflect pre-training rather than post-training.

% ..............................................
\subsection{Reproduction Models}
\label{subsec:reproduction-models}

To validate our implementation against \citet{kargaran2025mexa}, we
re-evaluate two models with published reference scores:
\textbf{Llama~3.1~8B} \citep{dubey2024llama3}, the primary English-centric
decoder studied in the original paper, and \textbf{Mistral~7B~v0.3}
\citep{jiang2023mistral}. Agreement with the published
$\mu_{\text{Max}}$/$\mu_{\text{Mean}}$ values on the Table-1
configurations (Section~\ref{sec:datasets}) establishes that our
embedding-extraction and scoring code reproduces the original pipeline,
and anchors all subsequent extensions.

% ..............................................
\subsection{New Decoder Families}
\label{subsec:new-decoders}

The first extension applies MEXA to decoder families released after the
original study, covering different data mixtures and scales:

\begin{itemize}
    \item \textbf{Qwen3} \citep{yang2025qwen3} at 0.6B, 1.7B, 4B, and 8B
    parameters --- a family trained with a substantially more multilingual
    data mixture than Llama, whose within-family size sweep isolates the
    effect of scale under a fixed recipe;
    \item \textbf{Qwen3.5~9B}, the successor generation, to measure
    recipe-level progress at approximately constant parameter count;
    \item \textbf{Apertus~8B} and \textbf{Apertus-Mini~4B}, openly trained
    models with an explicitly multilingual, compliance-focused data
    policy, providing a counterpoint to English-centric training;
    \item \textbf{Gemma~4} in five variants (E2B, E4B, 12B, 31B, and the
    sparse 26B-A4B), Google's current open-weight family, spanning
    efficient edge-scale to 31B-parameter models.
\end{itemize}

\noindent
Together with the reproduction models, this yields size ladders within
families (Qwen3 0.6B$\to$8B, Gemma~4 E2B$\to$31B) and recipe contrasts
across families at matched scale (Llama~3.1~8B vs.\ Qwen3~8B vs.\
Apertus~8B).

% ..............................................
\subsection{Mixture-of-Experts Models}
\label{subsec:moe-models}

Sparse mixture-of-experts (MoE) architectures route each token through a
small subset of expert feed-forward networks, decoupling parameter count
from per-token compute \citep{jiang2024mixtral}. Whether expert routing
helps or hurts cross-lingual alignment is an open question: experts could
specialize by language (fragmenting the shared space) or by function
(leaving alignment intact). We evaluate \textbf{Mixtral~8x7B} and
\textbf{Mixtral~8x22B} \citep{jiang2024mixtral}, comparing them against
their dense sibling Mistral~7B~v0.3, and \textbf{Gemma~4~26B-A4B},
comparing it against dense Gemma~4 variants. In all cases the MEXA
pipeline applies unchanged, since it only consumes the post-layer hidden
states.

% ..............................................
\subsection{Encoder and Embedding Models}
\label{subsec:encoder-models}

The final model axis contrasts causal decoders with bidirectional
encoders and purpose-built sentence-embedding models, which provide an
upper-bound reference for what alignment looks like when it is an explicit
training objective rather than a by-product:

\begin{itemize}
    \item \textbf{Masked-LM encoders:} XLM-R base and large
    \citep{conneau2020xlmr}, Glot500 \citep{imani2023glot500} (an XLM-R
    variant extended to 500+ languages, particularly relevant for the
    Bible corpus tail), and mmBERT-base;
    \item \textbf{Contrastively trained sentence encoders:} LaBSE
    \citep{feng2022labse} and multilingual-E5-base
    \citep{wang2024multilingual}, both explicitly optimized so that
    translations are nearest neighbors;
    \item \textbf{Decoder-based embedding models:} Qwen3-Embedding at
    0.6B, 4B, and 8B, which share the Qwen3 backbone but are fine-tuned
    for retrieval --- allowing a controlled comparison of the same
    architecture before and after embedding-specific training.
\end{itemize}

% ------------------------------------------------------------
\section{Datasets}
\label{sec:datasets}

All experiments draw on the two parallel corpora introduced in
Section~\ref{sec:parallel-corpora}, in configurations chosen to match the
original paper where reproduction is intended and to extend coverage
elsewhere.

\paragraph{FLORES-200.}
From the FLORES-200 devtest split \citep{nllbteam2022} we use:
\begin{itemize}
    \item \textbf{FLORES Table-1 subset:} the 116 languages overlapping
    with Belebele, with $n = 100$ parallel sentences per language ---
    the exact configuration of Table~1 in \citet{kargaran2025mexa}, used
    for reproduction and for all model-family comparisons;
    \item \textbf{FLORES Table-1 (2000):} the same 116 languages with all
    $n = 2{,}000$ devtest sentences, to quantify the effect of sample
    size on the retrieval criterion (larger $n$ means more distractors
    and therefore a strictly harder task);
    \item \textbf{FLORES Full:} all 204 language--script pairs, for
    maximal high-quality coverage.
\end{itemize}

\paragraph{Parallel Bible Corpus (sPBC).}
From the super-parallel Bible corpus \citep{mayer2014bible} we use:
\begin{itemize}
    \item \textbf{Bible Table-1 subset:} 101 languages with $n = 103$
    parallel verses, matching the original paper's Bible configuration;
    \item \textbf{Bible Full:} the complete corpus of more than 1{,}400
    language--script pairs, which extends evaluation far into the
    low-resource tail where no downstream benchmark exists.
\end{itemize}

English (\texttt{eng\_Latn}) serves as the pivot in all standard runs; the
pivot-substitution runs of Section~\ref{sec:mexa-pipeline} use the FLORES
Table-1 subset. Language identifiers follow the FLORES-200 convention of
ISO~639-3 code plus script (e.g.\ \texttt{shn\_Mymr}).

% ------------------------------------------------------------
\section{Experimental Settings}
\label{sec:experimental-settings}

% ..............................................
\subsection{Evaluation and Analysis Methods}
\label{subsec:analysis-methods}

Unless stated otherwise, reported scores use position-weighted average
embeddings (Equation~\eqref{eq:weighted-avg}) and are given as
$\mu_{\text{Max}}$ and $\mu_{\text{Mean}}$ per model--dataset
configuration; last-token embeddings and the centered variant are reported
as ablations. The analysis proceeds on four levels:

\begin{itemize}
    \item \textbf{Model level:} scores averaged over all languages of a
    configuration, enabling family- and scale-comparisons and the
    reproduction check against published values;
    \item \textbf{Language level:} per-language scores, ranked and grouped
    by script and by resource level, and cross-referenced with tokenizer
    fertility (computed on both corpora) to separate tokenization effects
    from representational effects;
    \item \textbf{Layer level:} full layer-wise alignment profiles
    $\mu C(\cdot, l)$, visualized as heatmaps, to locate where in the
    network alignment peaks and to test the intermediate-layer prediction
    of the pivot hypothesis;
    \item \textbf{Validity level:} Pearson correlations between MEXA
    scores and downstream accuracy on Belebele
    \citep{bandarkar2024belebele} and m-ARC \citep{lai2023okapi} over the
    languages both cover, testing whether alignment is a faithful proxy
    for task performance; cross-model Pearson correlations of per-language
    score profiles additionally measure how consistent alignment
    structure is across architectures.
\end{itemize}

For qualitative inspection, sentence embeddings are also projected to two
and three dimensions with PCA and t-SNE, colored by language and script,
to visualize the clustering behavior that the retrieval score summarizes
numerically.

% ..............................................
\subsection{Implementation and Compute}
\label{subsec:implementation}

The pipeline is implemented in Python with PyTorch and Hugging Face
Transformers. \texttt{embed\_extractor.py} loads each checkpoint via
\texttt{AutoModelForCausalLM} (or \texttt{AutoModel} for encoders) with
output of all hidden states enabled, tokenizes each sentence with the
model's own tokenizer, and stores per-layer sentence embeddings as
serialized NumPy arrays, one file per language.
\texttt{compute\_mexa.py} then builds the layer-wise cosine similarity
matrices in float64 and evaluates Equation~\eqref{eq:mexa}; because
scoring operates on cached embeddings, ablations over pooling strategies,
pivots, and centering re-use a single (expensive) extraction pass per
model--corpus pair. Formatted per-language and per-layer scores are
exported to CSV/JSON and explored in an interactive React dashboard
developed for this thesis.

Small models were run locally; all large-scale extraction jobs (8B and
above, MoE, and full-corpus sweeps) were executed on the CIT-TUM-HN
SLURM cluster at TUM Campus Heilbronn on NVIDIA data-center GPUs, with
models loaded in reduced precision (bfloat16, with quantized loading for
the largest MoE checkpoints) and embeddings cast to float32 before
storage. Random effects play no role in the pipeline --- extraction and
scoring are deterministic given a checkpoint and corpus --- so no seed
averaging is required.
 from a report/book-class document
% (e.g. the tum-thesis-latex template). Requires: amsmath, amssymb,
% booktabs, natbib, and the accompanying references.bib.
% ============================================================


% ============================================================
% CHAPTER 2 — BACKGROUND
% ============================================================
\chapter{Background}
\label{ch:background}

This chapter introduces the concepts on which the thesis builds. We first
review how multilingual capabilities arise in large language models (LLMs)
and where they break down (Section~\ref{sec:multilingual-lms}). We then
discuss the special role of English as a pivot language, both as an explicit
engineering device in classical multilingual NLP and as an implicit property
of the internal representations of English-centric LLMs
(Section~\ref{sec:english-pivot}). Section~\ref{sec:multilingual-eval}
surveys how multilingual ability is evaluated today and motivates
alignment-based, reference-free metrics such as MEXA. Finally,
Section~\ref{sec:parallel-corpora} describes the parallel corpora that make
such evaluations possible and the resource disparities that constrain them.

% ------------------------------------------------------------
\section{Multilingual Language Models}
\label{sec:multilingual-lms}

% ..............................................
\subsection{Architectures and Pre-training}
\label{subsec:architectures}

Virtually all modern language models are based on the Transformer
architecture \citep{vaswani2017attention}, a stack of identical layers that
combine multi-head self-attention with position-wise feed-forward networks.
Each layer transforms a sequence of token representations into a new
sequence of contextualized representations; the output of layer $l$ for
token $t$, the \emph{hidden state} $h_t^{(l)}$, is the central object of
study in this thesis, because it is from these layer-wise hidden states
that we extract sentence embeddings.

Two architectural families dominate multilingual NLP. \emph{Encoder-only}
models such as mBERT \citep{devlin2019bert} and XLM-R
\citep{conneau2020xlmr} use bidirectional attention and are pre-trained
with masked-language-modeling objectives on deliberately multilingual
corpora (XLM-R covers 100 languages sampled from CommonCrawl). Their
representations are designed to serve downstream classification and
retrieval tasks, and dedicated sentence encoders such as LaBSE
\citep{feng2022labse}, LASER3 \citep{heffernan2022laser3}, and multilingual
E5 \citep{wang2024multilingual} additionally fine-tune on parallel data
with contrastive or translation-ranking objectives so that translations map
to nearby points in embedding space. \emph{Decoder-only} models such as
GPT-3 \citep{brown2020gpt3} and the Llama family \citep{dubey2024llama3}
use causal (left-to-right) attention and are pre-trained autoregressively
to predict the next token. They are the dominant paradigm for generative
LLMs, but their pre-training corpora are heavily skewed toward English:
for Llama~3, roughly 90\% of the pre-training tokens are English
\citep{dubey2024llama3}. Models of this kind are commonly called
\emph{English-centric}, even though they exhibit non-trivial abilities in
dozens of languages.

This asymmetry creates the central question of the thesis: an
English-centric decoder model is never explicitly trained to align
languages, yet it acquires multilingual competence as a by-product of
incidental multilingual data in its corpus. Understanding \emph{where} and
\emph{how well} this competence materializes inside the network --- rather
than only measuring it behaviorally on downstream tasks --- is what
alignment-based evaluation aims to do.

% ..............................................
\subsection{The Tokenization Bottleneck}
\label{subsec:tokenization}

Before any text reaches the Transformer layers, it is segmented into
subword units by a tokenizer, typically trained with byte-pair encoding
\citep{sennrich2016bpe} or a unigram language model as implemented in
SentencePiece \citep{kudo2018sentencepiece}. The tokenizer's vocabulary is
learned on (a sample of) the pre-training corpus, and therefore inherits
its language distribution: frequent English words receive dedicated tokens,
while words in under-represented languages are shattered into many short
fragments, in the worst case into individual bytes.

The severity of this effect is commonly quantified by \emph{fertility},
the average number of subword tokens produced per word (or per character
for scripts without whitespace word boundaries). High fertility has three
practical consequences. First, it inflates sequence lengths, making
inference disproportionately expensive for low-resource languages. Second,
it dilutes the semantic content carried by each token: a model that sees
Amharic or Tibetan only as byte sequences must reconstruct lexical meaning
from much noisier units than it does for English. Third, and most relevant
to this thesis, it affects the quality of sentence embeddings built from
token hidden states, since averaging over many weakly informative fragments
yields blurrier sentence representations. Non-Latin scripts (Ge'ez,
Myanmar, Khmer, Cherokee) and low-resource languages regardless of script
are the most affected. Throughout our experiments we therefore track
tokenizer fertility on both evaluation corpora alongside the alignment
scores, which allows us to distinguish genuinely weak cross-lingual
representations from representations that are handicapped already at the
tokenization stage.

% ..............................................
\subsection{Cross-lingual Transfer}
\label{subsec:transfer}

The practical value of multilingual models rests on \emph{cross-lingual
transfer}: capabilities acquired predominantly in one language (in
practice, English) become available in other languages without
language-specific supervision. The standard explanation is that the model
learns a partially \emph{shared semantic space} in which translation
equivalents --- words or sentences with the same meaning --- receive
similar representations, so that task knowledge attached to a region of
that space is accessible from any language that maps into it
\citep{hammerl2024survey}.

Empirical evidence supports this picture with an important refinement:
the degree of sharing is \emph{layer-dependent}. Input embeddings and the
first layers remain largely language-specific, dominated by lexical and
orthographic features; representations become increasingly
language-neutral through the middle of the network; and the final layers
re-specialize toward the output vocabulary of the target language.
Interpretability studies on Llama-family models show precisely this
trajectory \citep{wendler2024llamas,zhao2024multilingualism}, and it
motivates the layer-wise design of MEXA: if cross-lingual alignment exists
anywhere in an English-centric model, it should be found in the
intermediate layers, and a single-layer (e.g., last-layer) probe would
systematically underestimate it.

\citet{hammerl2024survey} organize the methods that \emph{create} or
\emph{improve} such alignment into families --- joint multilingual
pre-training, alignment objectives on parallel data, contrastive
sentence-level training, and post-hoc mapping between monolingual spaces.
This thesis is concerned with the complementary problem: not producing
alignment, but \emph{measuring} the alignment that a given model already
has, layer by layer and language by language.

% ------------------------------------------------------------
\section{English as a Pivot Language}
\label{sec:english-pivot}

% ..............................................
\subsection{Pivot Translation and Alignment}
\label{subsec:pivot-translation}

Using English as a bridge between languages long predates neural LLMs. In
statistical machine translation, direct parallel data between two
non-English languages is often scarce, so translation is routed through a
\emph{pivot} language: source $\rightarrow$ English $\rightarrow$ target
\citep{wu2007pivot}. Because English participates in the largest number of
parallel corpora, it is the pivot of choice, and the quality of any
source--target system becomes bounded by the quality of the two
English-anchored legs. Neural machine translation reproduced the pattern
in a softer form: multilingual NMT systems trained on English-centric data
develop an ``interlingua-like'' internal representation that enables
zero-shot translation between language pairs never seen together in
training \citep{johnson2017zeroshot}.

Modern decoder-only LLMs internalize this bridge. Using the logit-lens
technique, \citet{wendler2024llamas} show that when Llama-2 models process
non-English input, the intermediate layers pass through a phase in which
the representations are closest to \emph{English} tokens before the final
layers rotate them toward the target-language vocabulary --- the model
demonstrably ``works in English'' in its middle layers.
\citet{zhao2024multilingualism} corroborate this for reasoning tasks,
finding that multilingual inputs are converted into English-centric latent
representations in which the actual computation takes place. This
\emph{pivot hypothesis} is the theoretical foundation of MEXA: if
non-English understanding is routed through English-like intermediate
representations, then the degree to which a language's representations
align with English is a direct measurement of how much of the model's
(English-acquired) capability that language can access.

% ..............................................
\subsection{English Centricity and Bias}
\label{subsec:english-bias}

The pivot mechanism is a double-edged sword. On one side, it is exactly
what makes cross-lingual transfer from a 90\%-English corpus possible at
all. On the other, it builds a structural asymmetry into the model.
Languages that map cleanly onto English-like representations inherit the
model's full capability; languages that do not --- because of typological
distance, script, or sheer data scarcity --- are effectively locked out,
regardless of how much task knowledge the model possesses. Concepts and
distinctions that English does not encode (grammatical evidentiality,
honorific systems, culture-specific terminology) risk being flattened when
they are forced through an English-shaped bottleneck, and cultural framing
in generated text tends to default to Anglophone norms.

For evaluation, English centricity raises a methodological question that
this thesis addresses empirically: is English \emph{privileged} as a
measurement anchor, or merely \emph{convenient}? MEXA, as proposed, always
measures alignment against English. If the pivot hypothesis is the true
mechanism, replacing English with another pivot language should
systematically lower the measured alignment, in proportion to how far the
alternative pivot is from the model's internal hub. Our non-English-pivot
experiments (Section~\ref{sec:experimental-settings}) test exactly this by
recomputing alignment against French, German, Arabic, Chinese, and Basque
pivots.

% ------------------------------------------------------------
\section{Multilingual Evaluation}
\label{sec:multilingual-eval}

% ..............................................
\subsection{Standard Benchmarks}
\label{subsec:benchmarks}

The most direct way to assess multilingual capability is to run the model
on parallel downstream tasks. FLORES-200 \citep{nllbteam2022} provides
professionally translated Wikimedia sentences in 204 language--script
pairs and is the de-facto standard for evaluating translation and
cross-lingual retrieval. Belebele \citep{bandarkar2024belebele} builds a
machine reading comprehension benchmark on top of FLORES passages,
covering 122 language variants with fully parallel multiple-choice
questions --- crucially, differences in scores across languages can only
stem from the language itself, not from differing content. Knowledge- and
reasoning-oriented English benchmarks such as MMLU
\citep{hendrycks2021mmlu} and ARC have been machine-translated into
26--31 languages to form m-MMLU and m-ARC \citep{lai2023okapi}.

These benchmarks are indispensable, but they have structural limits: they
exist for at most a couple of hundred languages (Belebele: 122, m-MMLU:
34, m-ARC: 31), they are expensive to extend because each new language
needs high-quality translated task data, and machine-translated variants
inherit translation noise. For the several thousand remaining languages,
task-based evaluation is simply unavailable --- which is the gap that
alignment-based estimation aims to fill. In this thesis, downstream
benchmarks play the role of \emph{validation targets}: MEXA scores are
correlated against Belebele and m-ARC performance to test whether
alignment is a faithful proxy for task capability.

% ..............................................
\subsection{Reference-Based vs.\ Reference-Free Metrics}
\label{subsec:metrics}

Evaluation metrics in multilingual NLP fall into two broad families.
\emph{Reference-based} metrics compare model output against gold
references. Surface-overlap metrics such as BLEU
\citep{papineni2002bleu} are cheap and language-agnostic in principle, but
they reward n-gram overlap rather than meaning, penalize legitimate
paraphrases, and interact badly with morphologically rich languages and
non-whitespace scripts. Learned metrics such as COMET \citep{rei2020comet}
correlate far better with human judgments by scoring
source--hypothesis--reference triples with a fine-tuned multilingual
encoder, but they require human references and a trained scoring model,
and their reliability degrades precisely for the low-resource languages
where evaluation is most needed.

\emph{Reference-free} (quality-estimation or model-based) metrics remove
the dependency on gold outputs. Embedding-similarity approaches compare
representations of source and translation directly
\citep{reimers2020multilingual}, and model-internal approaches --- the
category MEXA belongs to --- do not look at generated text at all.
Instead, they probe the model's hidden states on parallel \emph{inputs}
and quantify structure that is known to correlate with capability. The
advantages are substantial: no task data, no references, no additional
trained judge, and applicability to any language for which a hundred
parallel sentences exist. The corresponding risk is validity --- an
internal signal is only useful insofar as it predicts external behavior
--- which is why correlation with downstream benchmarks
(Section~\ref{subsec:benchmarks}) is an essential part of the methodology
rather than an afterthought.

% ..............................................
\subsection{Challenges in Cross-lingual Alignment Evaluation}
\label{subsec:alignment-challenges}

Measuring alignment from raw embedding geometry is less straightforward
than it appears, for reasons rooted in the geometry of learned
representation spaces.

\paragraph{Anisotropy.} Transformer embedding spaces are strongly
anisotropic: representations occupy a narrow cone rather than spreading
over the full space, so the cosine similarity between two \emph{unrelated}
sentences is already high \citep{rajaee2022isotropy}. In multilingual
models the effect is compounded by outlier dimensions with extreme
variance that dominate cosine computations; removing a handful of such
dimensions can change similarity rankings entirely
\citep{hammerl2023anisotropy}. Absolute cosine values are therefore not
comparable across models, layers, or languages.

\paragraph{Hubness.} In high-dimensional spaces, some points become
``hubs'' that appear among the nearest neighbors of disproportionately
many queries. A naive nearest-neighbor criterion over cosine similarities
rewards such hubs and distorts retrieval-based measurements.

\paragraph{Language identity vs.\ semantics.} Embeddings encode both
\emph{what} a sentence says and \emph{which language} it is written in.
Two sentences in the same language are often closer to each other than a
sentence and its translation, simply because of the shared language
signal. A useful alignment metric must isolate the semantic component
from this language-identity component.

MEXA's design responds to all three problems at once: instead of using
cosine similarity as an absolute quantity, it uses it only
\emph{comparatively}, asking whether the correct translation outranks all
distractors within the same similarity matrix (a bidirectional
precision-at-1 criterion, formalized in
Section~\ref{sec:mexa-pipeline}). Rank-based comparisons within a fixed
matrix are invariant to the global scale and offset distortions that
anisotropy introduces, and the bidirectional requirement suppresses hub
sentences that would pass a one-directional test. As an extension, this
thesis additionally evaluates a \emph{mean-centered} variant in which each
language's mean embedding is subtracted before computing similarities,
directly removing the language-identity component discussed above.

% ------------------------------------------------------------
\section{Parallel Corpora}
\label{sec:parallel-corpora}

% ..............................................
\subsection{Sources and Typology}
\label{subsec:corpora-sources}

Alignment-based evaluation requires \emph{parallel} data: the same
sentences translated across many languages, so that semantic content is
held constant while language varies. Parallel corpora differ along
several axes that matter for measurement quality: translation quality
(professional vs.\ crawled/mined, e.g.\ OPUS \citep{tiedemann2012opus}),
domain (news, web, religious text), degree of parallelism (bitexts vs.\
$n$-way parallel corpora), and language coverage.

Two corpora sit at opposite ends of the coverage--quality trade-off and
are both used in this thesis. \textbf{FLORES-200} \citep{nllbteam2022}
contains professionally translated, quality-assured sentences from
Wikimedia sources in 204 language--script pairs; it is $n$-way parallel,
modern-domain, and high quality, but its coverage stops at roughly 200
languages. The \textbf{Parallel Bible Corpus} \citep{mayer2014bible}
exploits the fact that the Bible is the most translated text in existence,
with verse-level alignment across more than 1{,}400 language--script
pairs; the derived \emph{super-parallel} subset (sPBC) restricts to verses
available in all covered languages. Its domain is narrow and archaic and
translation styles vary, but it is the only resource that reaches deep
into the low-resource tail, which is exactly where reference-free
evaluation is most valuable.

% ..............................................
\subsection{Resource Disparity}
\label{subsec:resource-disparity}

The distribution of parallel (and monolingual) data across the world's
languages is extremely skewed. A few dozen high-resource languages account
for the overwhelming majority of available text; thousands of languages
have effectively no digital footprint beyond a Bible translation. This
disparity propagates through the entire pipeline studied in this thesis:
low-resource languages receive less pre-training data (weaker
representations), fragmented tokenization
(Section~\ref{subsec:tokenization}), and no downstream benchmarks
(Section~\ref{subsec:benchmarks}) --- so their capability can be neither
learned as well nor measured as easily.

For evaluation design, resource disparity has two concrete consequences.
First, any claim about ``multilingual'' capability that is validated only
on high-resource languages is systematically optimistic; extending
measurement to the sPBC's low-resource tail is therefore a substantive
robustness test, not an ornament. Second, corpus effects must be
disentangled from language effects: a language may score lower on the
Bible corpus than on FLORES because the model is genuinely weaker in it,
or because archaic religious register is far from the model's training
distribution. Using both corpora on the same models, as done here, makes
this confound visible and quantifiable.


% ============================================================
% CHAPTER 3 — METHODOLOGY
% ============================================================
\chapter{Methodology}
\label{ch:methodology}

This chapter specifies the experimental  of the thesis. We first
formalize the MEXA pipeline as implemented in our codebase
(Section~\ref{sec:mexa-pipeline}), then describe the evaluated models
(Section~\ref{sec:models}), the parallel corpora and their subsets
(Section~\ref{sec:datasets}), and finally the experimental settings,
analysis methods, and compute environment
(Section~\ref{sec:experimental-settings}). The methodology follows
\citet{kargaran2025mexa} for the reproduction experiments and extends it
along four axes: new decoder families, mixture-of-experts architectures,
encoder and embedding models, and non-English pivot languages.

% ------------------------------------------------------------
\section{The MEXA Pipeline}
\label{sec:mexa-pipeline}

MEXA (Multilingual Evaluation via Cross-Lingual Alignment) estimates the
multilingual coverage of a language model $m$ from parallel sentences
alone, in four stages: sentence-embedding extraction, layer-wise
similarity computation, alignment scoring, and aggregation.

\paragraph{Stage 1: Layer-wise sentence embeddings.}
Let $\mathcal{S} = \{s_1, \dots, s_n\}$ be a set of $n$ parallel sentences
available in a target language $L_1$ and a pivot language $L_2$ (English,
unless stated otherwise). Each sentence is passed through the model, and
hidden states $h_t^{(l)}$ are collected for every token $t$ at every layer
$l \in \{1, \dots, L\}$, including the embedding layer. Token states are
reduced to a single sentence embedding $e^{(l)}$ per layer with one of two
strategies:

\begin{itemize}
    \item \textbf{Position-weighted average.} Because causal attention
    means early tokens see little context while late tokens see all of it,
    tokens are averaged with weights proportional to their position:
    \begin{equation}
        e^{(l)} \;=\; \sum_{t=1}^{T} w_t\, h_t^{(l)},
        \qquad
        w_t \;=\; \frac{t}{\sum_{k=1}^{T} k},
        \label{eq:weighted-avg}
    \end{equation}
    where $T$ is the number of non-padding tokens. This is the primary
    setting in all reported results.
    \item \textbf{Last token.} The hidden state of the final non-padding
    token, $e^{(l)} = h_T^{(l)}$, which under causal attention is the only
    position that has attended to the entire sentence. This setting is
    used as a robustness check.
\end{itemize}

For encoder-only models (Section~\ref{subsec:encoder-models}), whose
attention is bidirectional, the same two reductions are applied to the
encoder hidden states, making the pipeline architecture-agnostic.
Embeddings are stored per language and per layer, so that scoring is
decoupled from (and far cheaper than) extraction.

\paragraph{Stage 2: Similarity matrix.}
For each layer $l$, a cosine similarity matrix
$C(L_1, L_2, m, l) \in \mathbb{R}^{n \times n}$ is computed between the
$n$ target-language embeddings and the $n$ pivot-language embeddings:
$c_{ij}(l)$ is the cosine similarity between target sentence $i$ and pivot
sentence $j$. The diagonal entries $c_{ii}(l)$ correspond to true parallel
pairs; all off-diagonal entries are non-parallel distractors.

\paragraph{Stage 3: Alignment score.}
As discussed in Section~\ref{subsec:alignment-challenges}, absolute cosine
values are unreliable in anisotropic LLM embedding spaces. MEXA therefore
scores alignment with a \emph{bidirectional retrieval} (precision-at-1)
criterion: a sentence pair counts as aligned only if its true similarity
strictly exceeds every distractor in both its row and its column,

\begin{equation}
    \mu C(L_1, L_2, m, l)
    \;=\;
    \frac{1}{n}
    \sum_{i=1}^{n}
    \mathbb{1}\!\left(
        c_{ii}(l) \;>\; \max_{j \neq i}\,\{\, c_{ij}(l),\; c_{ji}(l) \,\}
    \right).
    \label{eq:mexa}
\end{equation}

The score lies in $[0, 1]$; the expected value under random embeddings is
approximately $1/n^2$ per pair (both the row and the column condition must
hold by chance), so with $n = 100$ sentences a random baseline scores
essentially zero, which validates the use of small parallel samples.

\paragraph{Stage 4: Aggregation across layers.}
The layer-wise scores are summarized into two model-level statistics:
\begin{equation}
    \mu_{\text{Max}}(L_1) = \max_{l} \; \mu C(L_1, L_2, m, l),
    \qquad
    \mu_{\text{Mean}}(L_1) = \frac{1}{L} \sum_{l=1}^{L} \mu C(L_1, L_2, m, l).
\end{equation}
$\mu_{\text{Max}}$ captures the peak alignment a model achieves anywhere
in its depth --- typically in the middle-to-late layers, consistent with
the pivot hypothesis --- while $\mu_{\text{Mean}}$ rewards models that
sustain alignment across many layers. Model-level scores reported in
Chapter~4 average these quantities over all evaluated languages.

\paragraph{Extension: mean-centered scoring.}
To separate semantic alignment from the language-identity signal
(Section~\ref{subsec:alignment-challenges}), we additionally compute a
\emph{centered} variant of Equation~\eqref{eq:mexa}: before building the
similarity matrix at layer $l$, each language's embeddings are shifted by
their own mean,
\begin{equation}
    \tilde{e}_i^{(l)} \;=\; e_i^{(l)} - \frac{1}{n}\sum_{k=1}^{n} e_k^{(l)},
\end{equation}
applied independently to the target and pivot sides. Centering removes the
dominant direction shared by all sentences of a language and typically
sharpens retrieval; comparing centered and uncentered scores quantifies
how much measured alignment is carried by language-mean offsets rather
than sentence-level semantics.

\paragraph{Extension: non-English pivots.}
Equation~\eqref{eq:mexa} does not privilege English mathematically; the
pivot $L_2$ is a free parameter. We exploit this to test the pivot
hypothesis directly, recomputing all alignment scores with French, German,
Arabic, Chinese, and Basque as pivot languages --- spanning high-resource
Latin-script, non-Latin-script, and a language isolate --- in both the
standard and centered variants.

% ------------------------------------------------------------
\section{Models}
\label{sec:models}

We evaluate 25+ publicly available checkpoints, grouped into four
categories that correspond to the four experimental axes of the thesis.
All models are obtained from the Hugging Face Hub and evaluated without
any fine-tuning; base (non-instruction-tuned) variants are used where
available so that scores reflect pre-training rather than post-training.

% ..............................................
\subsection{Reproduction Models}
\label{subsec:reproduction-models}

To validate our implementation against \citet{kargaran2025mexa}, we
re-evaluate two models with published reference scores:
\textbf{Llama~3.1~8B} \citep{dubey2024llama3}, the primary English-centric
decoder studied in the original paper, and \textbf{Mistral~7B~v0.3}
\citep{jiang2023mistral}. Agreement with the published
$\mu_{\text{Max}}$/$\mu_{\text{Mean}}$ values on the Table-1
configurations (Section~\ref{sec:datasets}) establishes that our
embedding-extraction and scoring code reproduces the original pipeline,
and anchors all subsequent extensions.

% ..............................................
\subsection{New Decoder Families}
\label{subsec:new-decoders}

The first extension applies MEXA to decoder families released after the
original study, covering different data mixtures and scales:

\begin{itemize}
    \item \textbf{Qwen3} \citep{yang2025qwen3} at 0.6B, 1.7B, 4B, and 8B
    parameters --- a family trained with a substantially more multilingual
    data mixture than Llama, whose within-family size sweep isolates the
    effect of scale under a fixed recipe;
    \item \textbf{Qwen3.5~9B}, the successor generation, to measure
    recipe-level progress at approximately constant parameter count;
    \item \textbf{Apertus~8B} and \textbf{Apertus-Mini~4B}, openly trained
    models with an explicitly multilingual, compliance-focused data
    policy, providing a counterpoint to English-centric training;
    \item \textbf{Gemma~4} in five variants (E2B, E4B, 12B, 31B, and the
    sparse 26B-A4B), Google's current open-weight family, spanning
    efficient edge-scale to 31B-parameter models.
\end{itemize}

\noindent
Together with the reproduction models, this yields size ladders within
families (Qwen3 0.6B$\to$8B, Gemma~4 E2B$\to$31B) and recipe contrasts
across families at matched scale (Llama~3.1~8B vs.\ Qwen3~8B vs.\
Apertus~8B).

% ..............................................
\subsection{Mixture-of-Experts Models}
\label{subsec:moe-models}

Sparse mixture-of-experts (MoE) architectures route each token through a
small subset of expert feed-forward networks, decoupling parameter count
from per-token compute \citep{jiang2024mixtral}. Whether expert routing
helps or hurts cross-lingual alignment is an open question: experts could
specialize by language (fragmenting the shared space) or by function
(leaving alignment intact). We evaluate \textbf{Mixtral~8x7B} and
\textbf{Mixtral~8x22B} \citep{jiang2024mixtral}, comparing them against
their dense sibling Mistral~7B~v0.3, and \textbf{Gemma~4~26B-A4B},
comparing it against dense Gemma~4 variants. In all cases the MEXA
pipeline applies unchanged, since it only consumes the post-layer hidden
states.

% ..............................................
\subsection{Encoder and Embedding Models}
\label{subsec:encoder-models}

The final model axis contrasts causal decoders with bidirectional
encoders and purpose-built sentence-embedding models, which provide an
upper-bound reference for what alignment looks like when it is an explicit
training objective rather than a by-product:

\begin{itemize}
    \item \textbf{Masked-LM encoders:} XLM-R base and large
    \citep{conneau2020xlmr}, Glot500 \citep{imani2023glot500} (an XLM-R
    variant extended to 500+ languages, particularly relevant for the
    Bible corpus tail), and mmBERT-base;
    \item \textbf{Contrastively trained sentence encoders:} LaBSE
    \citep{feng2022labse} and multilingual-E5-base
    \citep{wang2024multilingual}, both explicitly optimized so that
    translations are nearest neighbors;
    \item \textbf{Decoder-based embedding models:} Qwen3-Embedding at
    0.6B, 4B, and 8B, which share the Qwen3 backbone but are fine-tuned
    for retrieval --- allowing a controlled comparison of the same
    architecture before and after embedding-specific training.
\end{itemize}

% ------------------------------------------------------------
\section{Datasets}
\label{sec:datasets}

All experiments draw on the two parallel corpora introduced in
Section~\ref{sec:parallel-corpora}, in configurations chosen to match the
original paper where reproduction is intended and to extend coverage
elsewhere.

\paragraph{FLORES-200.}
From the FLORES-200 devtest split \citep{nllbteam2022} we use:
\begin{itemize}
    \item \textbf{FLORES Table-1 subset:} the 116 languages overlapping
    with Belebele, with $n = 100$ parallel sentences per language ---
    the exact configuration of Table~1 in \citet{kargaran2025mexa}, used
    for reproduction and for all model-family comparisons;
    \item \textbf{FLORES Table-1 (2000):} the same 116 languages with all
    $n = 2{,}000$ devtest sentences, to quantify the effect of sample
    size on the retrieval criterion (larger $n$ means more distractors
    and therefore a strictly harder task);
    \item \textbf{FLORES Full:} all 204 language--script pairs, for
    maximal high-quality coverage.
\end{itemize}

\paragraph{Parallel Bible Corpus (sPBC).}
From the super-parallel Bible corpus \citep{mayer2014bible} we use:
\begin{itemize}
    \item \textbf{Bible Table-1 subset:} 101 languages with $n = 103$
    parallel verses, matching the original paper's Bible configuration;
    \item \textbf{Bible Full:} the complete corpus of more than 1{,}400
    language--script pairs, which extends evaluation far into the
    low-resource tail where no downstream benchmark exists.
\end{itemize}

English (\texttt{eng\_Latn}) serves as the pivot in all standard runs; the
pivot-substitution runs of Section~\ref{sec:mexa-pipeline} use the FLORES
Table-1 subset. Language identifiers follow the FLORES-200 convention of
ISO~639-3 code plus script (e.g.\ \texttt{shn\_Mymr}).

% ------------------------------------------------------------
\section{Experimental Settings}
\label{sec:experimental-settings}

% ..............................................
\subsection{Evaluation and Analysis Methods}
\label{subsec:analysis-methods}

Unless stated otherwise, reported scores use position-weighted average
embeddings (Equation~\eqref{eq:weighted-avg}) and are given as
$\mu_{\text{Max}}$ and $\mu_{\text{Mean}}$ per model--dataset
configuration; last-token embeddings and the centered variant are reported
as ablations. The analysis proceeds on four levels:

\begin{itemize}
    \item \textbf{Model level:} scores averaged over all languages of a
    configuration, enabling family- and scale-comparisons and the
    reproduction check against published values;
    \item \textbf{Language level:} per-language scores, ranked and grouped
    by script and by resource level, and cross-referenced with tokenizer
    fertility (computed on both corpora) to separate tokenization effects
    from representational effects;
    \item \textbf{Layer level:} full layer-wise alignment profiles
    $\mu C(\cdot, l)$, visualized as heatmaps, to locate where in the
    network alignment peaks and to test the intermediate-layer prediction
    of the pivot hypothesis;
    \item \textbf{Validity level:} Pearson correlations between MEXA
    scores and downstream accuracy on Belebele
    \citep{bandarkar2024belebele} and m-ARC \citep{lai2023okapi} over the
    languages both cover, testing whether alignment is a faithful proxy
    for task performance; cross-model Pearson correlations of per-language
    score profiles additionally measure how consistent alignment
    structure is across architectures.
\end{itemize}

For qualitative inspection, sentence embeddings are also projected to two
and three dimensions with PCA and t-SNE, colored by language and script,
to visualize the clustering behavior that the retrieval score summarizes
numerically.

% ..............................................
\subsection{Implementation and Compute}
\label{subsec:implementation}

The pipeline is implemented in Python with PyTorch and Hugging Face
Transformers. \texttt{embed\_extractor.py} loads each checkpoint via
\texttt{AutoModelForCausalLM} (or \texttt{AutoModel} for encoders) with
output of all hidden states enabled, tokenizes each sentence with the
model's own tokenizer, and stores per-layer sentence embeddings as
serialized NumPy arrays, one file per language.
\texttt{compute\_mexa.py} then builds the layer-wise cosine similarity
matrices in float64 and evaluates Equation~\eqref{eq:mexa}; because
scoring operates on cached embeddings, ablations over pooling strategies,
pivots, and centering re-use a single (expensive) extraction pass per
model--corpus pair. Formatted per-language and per-layer scores are
exported to CSV/JSON and explored in an interactive React dashboard
developed for this thesis.

Small models were run locally; all large-scale extraction jobs (8B and
above, MoE, and full-corpus sweeps) were executed on the CIT-TUM-HN
SLURM cluster at TUM Campus Heilbronn on NVIDIA data-center GPUs, with
models loaded in reduced precision (bfloat16, with quantized loading for
the largest MoE checkpoints) and embeddings cast to float32 before
storage. Random effects play no role in the pipeline --- extraction and
scoring are deterministic given a checkpoint and corpus --- so no seed
averaging is required.
 from a report/book-class document
% (e.g. the tum-thesis-latex template). Requires: amsmath, amssymb,
% booktabs, natbib, and the accompanying references.bib.
% ============================================================


% ============================================================
% CHAPTER 2 — BACKGROUND
% ============================================================
\chapter{Background}
\label{ch:background}

This chapter introduces the concepts on which the thesis builds. We first
review how multilingual capabilities arise in large language models (LLMs)
and where they break down (Section~\ref{sec:multilingual-lms}). We then
discuss the special role of English as a pivot language, both as an explicit
engineering device in classical multilingual NLP and as an implicit property
of the internal representations of English-centric LLMs
(Section~\ref{sec:english-pivot}). Section~\ref{sec:multilingual-eval}
surveys how multilingual ability is evaluated today and motivates
alignment-based, reference-free metrics such as MEXA. Finally,
Section~\ref{sec:parallel-corpora} describes the parallel corpora that make
such evaluations possible and the resource disparities that constrain them.

% ------------------------------------------------------------
\section{Multilingual Language Models}
\label{sec:multilingual-lms}

% ..............................................
\subsection{Architectures and Pre-training}
\label{subsec:architectures}

Virtually all modern language models are based on the Transformer
architecture \citep{vaswani2017attention}, a stack of identical layers that
combine multi-head self-attention with position-wise feed-forward networks.
Each layer transforms a sequence of token representations into a new
sequence of contextualized representations; the output of layer $l$ for
token $t$, the \emph{hidden state} $h_t^{(l)}$, is the central object of
study in this thesis, because it is from these layer-wise hidden states
that we extract sentence embeddings.

Two architectural families dominate multilingual NLP. \emph{Encoder-only}
models such as mBERT \citep{devlin2019bert} and XLM-R
\citep{conneau2020xlmr} use bidirectional attention and are pre-trained
with masked-language-modeling objectives on deliberately multilingual
corpora (XLM-R covers 100 languages sampled from CommonCrawl). Their
representations are designed to serve downstream classification and
retrieval tasks, and dedicated sentence encoders such as LaBSE
\citep{feng2022labse}, LASER3 \citep{heffernan2022laser3}, and multilingual
E5 \citep{wang2024multilingual} additionally fine-tune on parallel data
with contrastive or translation-ranking objectives so that translations map
to nearby points in embedding space. \emph{Decoder-only} models such as
GPT-3 \citep{brown2020gpt3} and the Llama family \citep{dubey2024llama3}
use causal (left-to-right) attention and are pre-trained autoregressively
to predict the next token. They are the dominant paradigm for generative
LLMs, but their pre-training corpora are heavily skewed toward English:
for Llama~3, roughly 90\% of the pre-training tokens are English
\citep{dubey2024llama3}. Models of this kind are commonly called
\emph{English-centric}, even though they exhibit non-trivial abilities in
dozens of languages.

This asymmetry creates the central question of the thesis: an
English-centric decoder model is never explicitly trained to align
languages, yet it acquires multilingual competence as a by-product of
incidental multilingual data in its corpus. Understanding \emph{where} and
\emph{how well} this competence materializes inside the network --- rather
than only measuring it behaviorally on downstream tasks --- is what
alignment-based evaluation aims to do.

% ..............................................
\subsection{The Tokenization Bottleneck}
\label{subsec:tokenization}

Before any text reaches the Transformer layers, it is segmented into
subword units by a tokenizer, typically trained with byte-pair encoding
\citep{sennrich2016bpe} or a unigram language model as implemented in
SentencePiece \citep{kudo2018sentencepiece}. The tokenizer's vocabulary is
learned on (a sample of) the pre-training corpus, and therefore inherits
its language distribution: frequent English words receive dedicated tokens,
while words in under-represented languages are shattered into many short
fragments, in the worst case into individual bytes.

The severity of this effect is commonly quantified by \emph{fertility},
the average number of subword tokens produced per word (or per character
for scripts without whitespace word boundaries). High fertility has three
practical consequences. First, it inflates sequence lengths, making
inference disproportionately expensive for low-resource languages. Second,
it dilutes the semantic content carried by each token: a model that sees
Amharic or Tibetan only as byte sequences must reconstruct lexical meaning
from much noisier units than it does for English. Third, and most relevant
to this thesis, it affects the quality of sentence embeddings built from
token hidden states, since averaging over many weakly informative fragments
yields blurrier sentence representations. Non-Latin scripts (Ge'ez,
Myanmar, Khmer, Cherokee) and low-resource languages regardless of script
are the most affected. Throughout our experiments we therefore track
tokenizer fertility on both evaluation corpora alongside the alignment
scores, which allows us to distinguish genuinely weak cross-lingual
representations from representations that are handicapped already at the
tokenization stage.

% ..............................................
\subsection{Cross-lingual Transfer}
\label{subsec:transfer}

The practical value of multilingual models rests on \emph{cross-lingual
transfer}: capabilities acquired predominantly in one language (in
practice, English) become available in other languages without
language-specific supervision. The standard explanation is that the model
learns a partially \emph{shared semantic space} in which translation
equivalents --- words or sentences with the same meaning --- receive
similar representations, so that task knowledge attached to a region of
that space is accessible from any language that maps into it
\citep{hammerl2024survey}.

Empirical evidence supports this picture with an important refinement:
the degree of sharing is \emph{layer-dependent}. Input embeddings and the
first layers remain largely language-specific, dominated by lexical and
orthographic features; representations become increasingly
language-neutral through the middle of the network; and the final layers
re-specialize toward the output vocabulary of the target language.
Interpretability studies on Llama-family models show precisely this
trajectory \citep{wendler2024llamas,zhao2024multilingualism}, and it
motivates the layer-wise design of MEXA: if cross-lingual alignment exists
anywhere in an English-centric model, it should be found in the
intermediate layers, and a single-layer (e.g., last-layer) probe would
systematically underestimate it.

\citet{hammerl2024survey} organize the methods that \emph{create} or
\emph{improve} such alignment into families --- joint multilingual
pre-training, alignment objectives on parallel data, contrastive
sentence-level training, and post-hoc mapping between monolingual spaces.
This thesis is concerned with the complementary problem: not producing
alignment, but \emph{measuring} the alignment that a given model already
has, layer by layer and language by language.

% ------------------------------------------------------------
\section{English as a Pivot Language}
\label{sec:english-pivot}

% ..............................................
\subsection{Pivot Translation and Alignment}
\label{subsec:pivot-translation}

Using English as a bridge between languages long predates neural LLMs. In
statistical machine translation, direct parallel data between two
non-English languages is often scarce, so translation is routed through a
\emph{pivot} language: source $\rightarrow$ English $\rightarrow$ target
\citep{wu2007pivot}. Because English participates in the largest number of
parallel corpora, it is the pivot of choice, and the quality of any
source--target system becomes bounded by the quality of the two
English-anchored legs. Neural machine translation reproduced the pattern
in a softer form: multilingual NMT systems trained on English-centric data
develop an ``interlingua-like'' internal representation that enables
zero-shot translation between language pairs never seen together in
training \citep{johnson2017zeroshot}.

Modern decoder-only LLMs internalize this bridge. Using the logit-lens
technique, \citet{wendler2024llamas} show that when Llama-2 models process
non-English input, the intermediate layers pass through a phase in which
the representations are closest to \emph{English} tokens before the final
layers rotate them toward the target-language vocabulary --- the model
demonstrably ``works in English'' in its middle layers.
\citet{zhao2024multilingualism} corroborate this for reasoning tasks,
finding that multilingual inputs are converted into English-centric latent
representations in which the actual computation takes place. This
\emph{pivot hypothesis} is the theoretical foundation of MEXA: if
non-English understanding is routed through English-like intermediate
representations, then the degree to which a language's representations
align with English is a direct measurement of how much of the model's
(English-acquired) capability that language can access.

% ..............................................
\subsection{English Centricity and Bias}
\label{subsec:english-bias}

The pivot mechanism is a double-edged sword. On one side, it is exactly
what makes cross-lingual transfer from a 90\%-English corpus possible at
all. On the other, it builds a structural asymmetry into the model.
Languages that map cleanly onto English-like representations inherit the
model's full capability; languages that do not --- because of typological
distance, script, or sheer data scarcity --- are effectively locked out,
regardless of how much task knowledge the model possesses. Concepts and
distinctions that English does not encode (grammatical evidentiality,
honorific systems, culture-specific terminology) risk being flattened when
they are forced through an English-shaped bottleneck, and cultural framing
in generated text tends to default to Anglophone norms.

For evaluation, English centricity raises a methodological question that
this thesis addresses empirically: is English \emph{privileged} as a
measurement anchor, or merely \emph{convenient}? MEXA, as proposed, always
measures alignment against English. If the pivot hypothesis is the true
mechanism, replacing English with another pivot language should
systematically lower the measured alignment, in proportion to how far the
alternative pivot is from the model's internal hub. Our non-English-pivot
experiments (Section~\ref{sec:experimental-settings}) test exactly this by
recomputing alignment against French, German, Arabic, Chinese, and Basque
pivots.

% ------------------------------------------------------------
\section{Multilingual Evaluation}
\label{sec:multilingual-eval}

% ..............................................
\subsection{Standard Benchmarks}
\label{subsec:benchmarks}

The most direct way to assess multilingual capability is to run the model
on parallel downstream tasks. FLORES-200 \citep{nllbteam2022} provides
professionally translated Wikimedia sentences in 204 language--script
pairs and is the de-facto standard for evaluating translation and
cross-lingual retrieval. Belebele \citep{bandarkar2024belebele} builds a
machine reading comprehension benchmark on top of FLORES passages,
covering 122 language variants with fully parallel multiple-choice
questions --- crucially, differences in scores across languages can only
stem from the language itself, not from differing content. Knowledge- and
reasoning-oriented English benchmarks such as MMLU
\citep{hendrycks2021mmlu} and ARC have been machine-translated into
26--31 languages to form m-MMLU and m-ARC \citep{lai2023okapi}.

These benchmarks are indispensable, but they have structural limits: they
exist for at most a couple of hundred languages (Belebele: 122, m-MMLU:
34, m-ARC: 31), they are expensive to extend because each new language
needs high-quality translated task data, and machine-translated variants
inherit translation noise. For the several thousand remaining languages,
task-based evaluation is simply unavailable --- which is the gap that
alignment-based estimation aims to fill. In this thesis, downstream
benchmarks play the role of \emph{validation targets}: MEXA scores are
correlated against Belebele and m-ARC performance to test whether
alignment is a faithful proxy for task capability.

% ..............................................
\subsection{Reference-Based vs.\ Reference-Free Metrics}
\label{subsec:metrics}

Evaluation metrics in multilingual NLP fall into two broad families.
\emph{Reference-based} metrics compare model output against gold
references. Surface-overlap metrics such as BLEU
\citep{papineni2002bleu} are cheap and language-agnostic in principle, but
they reward n-gram overlap rather than meaning, penalize legitimate
paraphrases, and interact badly with morphologically rich languages and
non-whitespace scripts. Learned metrics such as COMET \citep{rei2020comet}
correlate far better with human judgments by scoring
source--hypothesis--reference triples with a fine-tuned multilingual
encoder, but they require human references and a trained scoring model,
and their reliability degrades precisely for the low-resource languages
where evaluation is most needed.

\emph{Reference-free} (quality-estimation or model-based) metrics remove
the dependency on gold outputs. Embedding-similarity approaches compare
representations of source and translation directly
\citep{reimers2020multilingual}, and model-internal approaches --- the
category MEXA belongs to --- do not look at generated text at all.
Instead, they probe the model's hidden states on parallel \emph{inputs}
and quantify structure that is known to correlate with capability. The
advantages are substantial: no task data, no references, no additional
trained judge, and applicability to any language for which a hundred
parallel sentences exist. The corresponding risk is validity --- an
internal signal is only useful insofar as it predicts external behavior
--- which is why correlation with downstream benchmarks
(Section~\ref{subsec:benchmarks}) is an essential part of the methodology
rather than an afterthought.

% ..............................................
\subsection{Challenges in Cross-lingual Alignment Evaluation}
\label{subsec:alignment-challenges}

Measuring alignment from raw embedding geometry is less straightforward
than it appears, for reasons rooted in the geometry of learned
representation spaces.

\paragraph{Anisotropy.} Transformer embedding spaces are strongly
anisotropic: representations occupy a narrow cone rather than spreading
over the full space, so the cosine similarity between two \emph{unrelated}
sentences is already high \citep{rajaee2022isotropy}. In multilingual
models the effect is compounded by outlier dimensions with extreme
variance that dominate cosine computations; removing a handful of such
dimensions can change similarity rankings entirely
\citep{hammerl2023anisotropy}. Absolute cosine values are therefore not
comparable across models, layers, or languages.

\paragraph{Hubness.} In high-dimensional spaces, some points become
``hubs'' that appear among the nearest neighbors of disproportionately
many queries. A naive nearest-neighbor criterion over cosine similarities
rewards such hubs and distorts retrieval-based measurements.

\paragraph{Language identity vs.\ semantics.} Embeddings encode both
\emph{what} a sentence says and \emph{which language} it is written in.
Two sentences in the same language are often closer to each other than a
sentence and its translation, simply because of the shared language
signal. A useful alignment metric must isolate the semantic component
from this language-identity component.

MEXA's design responds to all three problems at once: instead of using
cosine similarity as an absolute quantity, it uses it only
\emph{comparatively}, asking whether the correct translation outranks all
distractors within the same similarity matrix (a bidirectional
precision-at-1 criterion, formalized in
Section~\ref{sec:mexa-pipeline}). Rank-based comparisons within a fixed
matrix are invariant to the global scale and offset distortions that
anisotropy introduces, and the bidirectional requirement suppresses hub
sentences that would pass a one-directional test. As an extension, this
thesis additionally evaluates a \emph{mean-centered} variant in which each
language's mean embedding is subtracted before computing similarities,
directly removing the language-identity component discussed above.

% ------------------------------------------------------------
\section{Parallel Corpora}
\label{sec:parallel-corpora}

% ..............................................
\subsection{Sources and Typology}
\label{subsec:corpora-sources}

Alignment-based evaluation requires \emph{parallel} data: the same
sentences translated across many languages, so that semantic content is
held constant while language varies. Parallel corpora differ along
several axes that matter for measurement quality: translation quality
(professional vs.\ crawled/mined, e.g.\ OPUS \citep{tiedemann2012opus}),
domain (news, web, religious text), degree of parallelism (bitexts vs.\
$n$-way parallel corpora), and language coverage.

Two corpora sit at opposite ends of the coverage--quality trade-off and
are both used in this thesis. \textbf{FLORES-200} \citep{nllbteam2022}
contains professionally translated, quality-assured sentences from
Wikimedia sources in 204 language--script pairs; it is $n$-way parallel,
modern-domain, and high quality, but its coverage stops at roughly 200
languages. The \textbf{Parallel Bible Corpus} \citep{mayer2014bible}
exploits the fact that the Bible is the most translated text in existence,
with verse-level alignment across more than 1{,}400 language--script
pairs; the derived \emph{super-parallel} subset (sPBC) restricts to verses
available in all covered languages. Its domain is narrow and archaic and
translation styles vary, but it is the only resource that reaches deep
into the low-resource tail, which is exactly where reference-free
evaluation is most valuable.

% ..............................................
\subsection{Resource Disparity}
\label{subsec:resource-disparity}

The distribution of parallel (and monolingual) data across the world's
languages is extremely skewed. A few dozen high-resource languages account
for the overwhelming majority of available text; thousands of languages
have effectively no digital footprint beyond a Bible translation. This
disparity propagates through the entire pipeline studied in this thesis:
low-resource languages receive less pre-training data (weaker
representations), fragmented tokenization
(Section~\ref{subsec:tokenization}), and no downstream benchmarks
(Section~\ref{subsec:benchmarks}) --- so their capability can be neither
learned as well nor measured as easily.

For evaluation design, resource disparity has two concrete consequences.
First, any claim about ``multilingual'' capability that is validated only
on high-resource languages is systematically optimistic; extending
measurement to the sPBC's low-resource tail is therefore a substantive
robustness test, not an ornament. Second, corpus effects must be
disentangled from language effects: a language may score lower on the
Bible corpus than on FLORES because the model is genuinely weaker in it,
or because archaic religious register is far from the model's training
distribution. Using both corpora on the same models, as done here, makes
this confound visible and quantifiable.


% ============================================================
% CHAPTER 3 — METHODOLOGY
% ============================================================
\chapter{Methodology}
\label{ch:methodology}

This chapter specifies the experimental  of the thesis. We first
formalize the MEXA pipeline as implemented in our codebase
(Section~\ref{sec:mexa-pipeline}), then describe the evaluated models
(Section~\ref{sec:models}), the parallel corpora and their subsets
(Section~\ref{sec:datasets}), and finally the experimental settings,
analysis methods, and compute environment
(Section~\ref{sec:experimental-settings}). The methodology follows
\citet{kargaran2025mexa} for the reproduction experiments and extends it
along four axes: new decoder families, mixture-of-experts architectures,
encoder and embedding models, and non-English pivot languages.

% ------------------------------------------------------------
\section{The MEXA Pipeline}
\label{sec:mexa-pipeline}

MEXA (Multilingual Evaluation via Cross-Lingual Alignment) estimates the
multilingual coverage of a language model $m$ from parallel sentences
alone, in four stages: sentence-embedding extraction, layer-wise
similarity computation, alignment scoring, and aggregation.

\paragraph{Stage 1: Layer-wise sentence embeddings.}
Let $\mathcal{S} = \{s_1, \dots, s_n\}$ be a set of $n$ parallel sentences
available in a target language $L_1$ and a pivot language $L_2$ (English,
unless stated otherwise). Each sentence is passed through the model, and
hidden states $h_t^{(l)}$ are collected for every token $t$ at every layer
$l \in \{1, \dots, L\}$, including the embedding layer. Token states are
reduced to a single sentence embedding $e^{(l)}$ per layer with one of two
strategies:

\begin{itemize}
    \item \textbf{Position-weighted average.} Because causal attention
    means early tokens see little context while late tokens see all of it,
    tokens are averaged with weights proportional to their position:
    \begin{equation}
        e^{(l)} \;=\; \sum_{t=1}^{T} w_t\, h_t^{(l)},
        \qquad
        w_t \;=\; \frac{t}{\sum_{k=1}^{T} k},
        \label{eq:weighted-avg}
    \end{equation}
    where $T$ is the number of non-padding tokens. This is the primary
    setting in all reported results.
    \item \textbf{Last token.} The hidden state of the final non-padding
    token, $e^{(l)} = h_T^{(l)}$, which under causal attention is the only
    position that has attended to the entire sentence. This setting is
    used as a robustness check.
\end{itemize}

For encoder-only models (Section~\ref{subsec:encoder-models}), whose
attention is bidirectional, the same two reductions are applied to the
encoder hidden states, making the pipeline architecture-agnostic.
Embeddings are stored per language and per layer, so that scoring is
decoupled from (and far cheaper than) extraction.

\paragraph{Stage 2: Similarity matrix.}
For each layer $l$, a cosine similarity matrix
$C(L_1, L_2, m, l) \in \mathbb{R}^{n \times n}$ is computed between the
$n$ target-language embeddings and the $n$ pivot-language embeddings:
$c_{ij}(l)$ is the cosine similarity between target sentence $i$ and pivot
sentence $j$. The diagonal entries $c_{ii}(l)$ correspond to true parallel
pairs; all off-diagonal entries are non-parallel distractors.

\paragraph{Stage 3: Alignment score.}
As discussed in Section~\ref{subsec:alignment-challenges}, absolute cosine
values are unreliable in anisotropic LLM embedding spaces. MEXA therefore
scores alignment with a \emph{bidirectional retrieval} (precision-at-1)
criterion: a sentence pair counts as aligned only if its true similarity
strictly exceeds every distractor in both its row and its column,

\begin{equation}
    \mu C(L_1, L_2, m, l)
    \;=\;
    \frac{1}{n}
    \sum_{i=1}^{n}
    \mathbb{1}\!\left(
        c_{ii}(l) \;>\; \max_{j \neq i}\,\{\, c_{ij}(l),\; c_{ji}(l) \,\}
    \right).
    \label{eq:mexa}
\end{equation}

The score lies in $[0, 1]$; the expected value under random embeddings is
approximately $1/n^2$ per pair (both the row and the column condition must
hold by chance), so with $n = 100$ sentences a random baseline scores
essentially zero, which validates the use of small parallel samples.

\paragraph{Stage 4: Aggregation across layers.}
The layer-wise scores are summarized into two model-level statistics:
\begin{equation}
    \mu_{\text{Max}}(L_1) = \max_{l} \; \mu C(L_1, L_2, m, l),
    \qquad
    \mu_{\text{Mean}}(L_1) = \frac{1}{L} \sum_{l=1}^{L} \mu C(L_1, L_2, m, l).
\end{equation}
$\mu_{\text{Max}}$ captures the peak alignment a model achieves anywhere
in its depth --- typically in the middle-to-late layers, consistent with
the pivot hypothesis --- while $\mu_{\text{Mean}}$ rewards models that
sustain alignment across many layers. Model-level scores reported in
Chapter~4 average these quantities over all evaluated languages.

\paragraph{Extension: mean-centered scoring.}
To separate semantic alignment from the language-identity signal
(Section~\ref{subsec:alignment-challenges}), we additionally compute a
\emph{centered} variant of Equation~\eqref{eq:mexa}: before building the
similarity matrix at layer $l$, each language's embeddings are shifted by
their own mean,
\begin{equation}
    \tilde{e}_i^{(l)} \;=\; e_i^{(l)} - \frac{1}{n}\sum_{k=1}^{n} e_k^{(l)},
\end{equation}
applied independently to the target and pivot sides. Centering removes the
dominant direction shared by all sentences of a language and typically
sharpens retrieval; comparing centered and uncentered scores quantifies
how much measured alignment is carried by language-mean offsets rather
than sentence-level semantics.

\paragraph{Extension: non-English pivots.}
Equation~\eqref{eq:mexa} does not privilege English mathematically; the
pivot $L_2$ is a free parameter. We exploit this to test the pivot
hypothesis directly, recomputing all alignment scores with French, German,
Arabic, Chinese, and Basque as pivot languages --- spanning high-resource
Latin-script, non-Latin-script, and a language isolate --- in both the
standard and centered variants.

% ------------------------------------------------------------
\section{Models}
\label{sec:models}

We evaluate 25+ publicly available checkpoints, grouped into four
categories that correspond to the four experimental axes of the thesis.
All models are obtained from the Hugging Face Hub and evaluated without
any fine-tuning; base (non-instruction-tuned) variants are used where
available so that scores reflect pre-training rather than post-training.

% ..............................................
\subsection{Reproduction Models}
\label{subsec:reproduction-models}

To validate our implementation against \citet{kargaran2025mexa}, we
re-evaluate two models with published reference scores:
\textbf{Llama~3.1~8B} \citep{dubey2024llama3}, the primary English-centric
decoder studied in the original paper, and \textbf{Mistral~7B~v0.3}
\citep{jiang2023mistral}. Agreement with the published
$\mu_{\text{Max}}$/$\mu_{\text{Mean}}$ values on the Table-1
configurations (Section~\ref{sec:datasets}) establishes that our
embedding-extraction and scoring code reproduces the original pipeline,
and anchors all subsequent extensions.

% ..............................................
\subsection{New Decoder Families}
\label{subsec:new-decoders}

The first extension applies MEXA to decoder families released after the
original study, covering different data mixtures and scales:

\begin{itemize}
    \item \textbf{Qwen3} \citep{yang2025qwen3} at 0.6B, 1.7B, 4B, and 8B
    parameters --- a family trained with a substantially more multilingual
    data mixture than Llama, whose within-family size sweep isolates the
    effect of scale under a fixed recipe;
    \item \textbf{Qwen3.5~9B}, the successor generation, to measure
    recipe-level progress at approximately constant parameter count;
    \item \textbf{Apertus~8B} and \textbf{Apertus-Mini~4B}, openly trained
    models with an explicitly multilingual, compliance-focused data
    policy, providing a counterpoint to English-centric training;
    \item \textbf{Gemma~4} in five variants (E2B, E4B, 12B, 31B, and the
    sparse 26B-A4B), Google's current open-weight family, spanning
    efficient edge-scale to 31B-parameter models.
\end{itemize}

\noindent
Together with the reproduction models, this yields size ladders within
families (Qwen3 0.6B$\to$8B, Gemma~4 E2B$\to$31B) and recipe contrasts
across families at matched scale (Llama~3.1~8B vs.\ Qwen3~8B vs.\
Apertus~8B).

% ..............................................
\subsection{Mixture-of-Experts Models}
\label{subsec:moe-models}

Sparse mixture-of-experts (MoE) architectures route each token through a
small subset of expert feed-forward networks, decoupling parameter count
from per-token compute \citep{jiang2024mixtral}. Whether expert routing
helps or hurts cross-lingual alignment is an open question: experts could
specialize by language (fragmenting the shared space) or by function
(leaving alignment intact). We evaluate \textbf{Mixtral~8x7B} and
\textbf{Mixtral~8x22B} \citep{jiang2024mixtral}, comparing them against
their dense sibling Mistral~7B~v0.3, and \textbf{Gemma~4~26B-A4B},
comparing it against dense Gemma~4 variants. In all cases the MEXA
pipeline applies unchanged, since it only consumes the post-layer hidden
states.

% ..............................................
\subsection{Encoder and Embedding Models}
\label{subsec:encoder-models}

The final model axis contrasts causal decoders with bidirectional
encoders and purpose-built sentence-embedding models, which provide an
upper-bound reference for what alignment looks like when it is an explicit
training objective rather than a by-product:

\begin{itemize}
    \item \textbf{Masked-LM encoders:} XLM-R base and large
    \citep{conneau2020xlmr}, Glot500 \citep{imani2023glot500} (an XLM-R
    variant extended to 500+ languages, particularly relevant for the
    Bible corpus tail), and mmBERT-base;
    \item \textbf{Contrastively trained sentence encoders:} LaBSE
    \citep{feng2022labse} and multilingual-E5-base
    \citep{wang2024multilingual}, both explicitly optimized so that
    translations are nearest neighbors;
    \item \textbf{Decoder-based embedding models:} Qwen3-Embedding at
    0.6B, 4B, and 8B, which share the Qwen3 backbone but are fine-tuned
    for retrieval --- allowing a controlled comparison of the same
    architecture before and after embedding-specific training.
\end{itemize}

% ------------------------------------------------------------
\section{Datasets}
\label{sec:datasets}

All experiments draw on the two parallel corpora introduced in
Section~\ref{sec:parallel-corpora}, in configurations chosen to match the
original paper where reproduction is intended and to extend coverage
elsewhere.

\paragraph{FLORES-200.}
From the FLORES-200 devtest split \citep{nllbteam2022} we use:
\begin{itemize}
    \item \textbf{FLORES Table-1 subset:} the 116 languages overlapping
    with Belebele, with $n = 100$ parallel sentences per language ---
    the exact configuration of Table~1 in \citet{kargaran2025mexa}, used
    for reproduction and for all model-family comparisons;
    \item \textbf{FLORES Table-1 (2000):} the same 116 languages with all
    $n = 2{,}000$ devtest sentences, to quantify the effect of sample
    size on the retrieval criterion (larger $n$ means more distractors
    and therefore a strictly harder task);
    \item \textbf{FLORES Full:} all 204 language--script pairs, for
    maximal high-quality coverage.
\end{itemize}

\paragraph{Parallel Bible Corpus (sPBC).}
From the super-parallel Bible corpus \citep{mayer2014bible} we use:
\begin{itemize}
    \item \textbf{Bible Table-1 subset:} 101 languages with $n = 103$
    parallel verses, matching the original paper's Bible configuration;
    \item \textbf{Bible Full:} the complete corpus of more than 1{,}400
    language--script pairs, which extends evaluation far into the
    low-resource tail where no downstream benchmark exists.
\end{itemize}

English (\texttt{eng\_Latn}) serves as the pivot in all standard runs; the
pivot-substitution runs of Section~\ref{sec:mexa-pipeline} use the FLORES
Table-1 subset. Language identifiers follow the FLORES-200 convention of
ISO~639-3 code plus script (e.g.\ \texttt{shn\_Mymr}).

% ------------------------------------------------------------
\section{Experimental Settings}
\label{sec:experimental-settings}

% ..............................................
\subsection{Evaluation and Analysis Methods}
\label{subsec:analysis-methods}

Unless stated otherwise, reported scores use position-weighted average
embeddings (Equation~\eqref{eq:weighted-avg}) and are given as
$\mu_{\text{Max}}$ and $\mu_{\text{Mean}}$ per model--dataset
configuration; last-token embeddings and the centered variant are reported
as ablations. The analysis proceeds on four levels:

\begin{itemize}
    \item \textbf{Model level:} scores averaged over all languages of a
    configuration, enabling family- and scale-comparisons and the
    reproduction check against published values;
    \item \textbf{Language level:} per-language scores, ranked and grouped
    by script and by resource level, and cross-referenced with tokenizer
    fertility (computed on both corpora) to separate tokenization effects
    from representational effects;
    \item \textbf{Layer level:} full layer-wise alignment profiles
    $\mu C(\cdot, l)$, visualized as heatmaps, to locate where in the
    network alignment peaks and to test the intermediate-layer prediction
    of the pivot hypothesis;
    \item \textbf{Validity level:} Pearson correlations between MEXA
    scores and downstream accuracy on Belebele
    \citep{bandarkar2024belebele} and m-ARC \citep{lai2023okapi} over the
    languages both cover, testing whether alignment is a faithful proxy
    for task performance; cross-model Pearson correlations of per-language
    score profiles additionally measure how consistent alignment
    structure is across architectures.
\end{itemize}

For qualitative inspection, sentence embeddings are also projected to two
and three dimensions with PCA and t-SNE, colored by language and script,
to visualize the clustering behavior that the retrieval score summarizes
numerically.

% ..............................................
\subsection{Implementation and Compute}
\label{subsec:implementation}

The pipeline is implemented in Python with PyTorch and Hugging Face
Transformers. \texttt{embed\_extractor.py} loads each checkpoint via
\texttt{AutoModelForCausalLM} (or \texttt{AutoModel} for encoders) with
output of all hidden states enabled, tokenizes each sentence with the
model's own tokenizer, and stores per-layer sentence embeddings as
serialized NumPy arrays, one file per language.
\texttt{compute\_mexa.py} then builds the layer-wise cosine similarity
matrices in float64 and evaluates Equation~\eqref{eq:mexa}; because
scoring operates on cached embeddings, ablations over pooling strategies,
pivots, and centering re-use a single (expensive) extraction pass per
model--corpus pair. Formatted per-language and per-layer scores are
exported to CSV/JSON and explored in an interactive React dashboard
developed for this thesis.

Small models were run locally; all large-scale extraction jobs (8B and
above, MoE, and full-corpus sweeps) were executed on the CIT-TUM-HN
SLURM cluster at TUM Campus Heilbronn on NVIDIA data-center GPUs, with
models loaded in reduced precision (bfloat16, with quantized loading for
the largest MoE checkpoints) and embeddings cast to float32 before
storage. Random effects play no role in the pipeline --- extraction and
scoring are deterministic given a checkpoint and corpus --- so no seed
averaging is required.
 from a report/book-class document
% (e.g. the tum-thesis-latex template). Requires: amsmath, amssymb,
% booktabs, natbib, and the accompanying references.bib.
% ============================================================


% ============================================================
% CHAPTER 2 — BACKGROUND
% ============================================================
\chapter{Background}
\label{ch:background}

This chapter introduces the concepts on which the thesis builds. We first
review how multilingual capabilities arise in large language models (LLMs)
and where they break down (Section~\ref{sec:multilingual-lms}). We then
discuss the special role of English as a pivot language, both as an explicit
engineering device in classical multilingual NLP and as an implicit property
of the internal representations of English-centric LLMs
(Section~\ref{sec:english-pivot}). Section~\ref{sec:multilingual-eval}
surveys how multilingual ability is evaluated today and motivates
alignment-based, reference-free metrics such as MEXA. Finally,
Section~\ref{sec:parallel-corpora} describes the parallel corpora that make
such evaluations possible and the resource disparities that constrain them.

% ------------------------------------------------------------
\section{Multilingual Language Models}
\label{sec:multilingual-lms}

% ..............................................
\subsection{Architectures and Pre-training}
\label{subsec:architectures}

Virtually all modern language models are based on the Transformer
architecture \citep{vaswani2017attention}, a stack of identical layers that
combine multi-head self-attention with position-wise feed-forward networks.
Each layer transforms a sequence of token representations into a new
sequence of contextualized representations; the output of layer $l$ for
token $t$, the \emph{hidden state} $h_t^{(l)}$, is the central object of
study in this thesis, because it is from these layer-wise hidden states
that we extract sentence embeddings.

Two architectural families dominate multilingual NLP. \emph{Encoder-only}
models such as mBERT \citep{devlin2019bert} and XLM-R
\citep{conneau2020xlmr} use bidirectional attention and are pre-trained
with masked-language-modeling objectives on deliberately multilingual
corpora (XLM-R covers 100 languages sampled from CommonCrawl). Their
representations are designed to serve downstream classification and
retrieval tasks, and dedicated sentence encoders such as LaBSE
\citep{feng2022labse}, LASER3 \citep{heffernan2022laser3}, and multilingual
E5 \citep{wang2024multilingual} additionally fine-tune on parallel data
with contrastive or translation-ranking objectives so that translations map
to nearby points in embedding space. \emph{Decoder-only} models such as
GPT-3 \citep{brown2020gpt3} and the Llama family \citep{dubey2024llama3}
use causal (left-to-right) attention and are pre-trained autoregressively
to predict the next token. They are the dominant paradigm for generative
LLMs, but their pre-training corpora are heavily skewed toward English:
for Llama~3, roughly 90\% of the pre-training tokens are English
\citep{dubey2024llama3}. Models of this kind are commonly called
\emph{English-centric}, even though they exhibit non-trivial abilities in
dozens of languages.

This asymmetry creates the central question of the thesis: an
English-centric decoder model is never explicitly trained to align
languages, yet it acquires multilingual competence as a by-product of
incidental multilingual data in its corpus. Understanding \emph{where} and
\emph{how well} this competence materializes inside the network --- rather
than only measuring it behaviorally on downstream tasks --- is what
alignment-based evaluation aims to do.

% ..............................................
\subsection{The Tokenization Bottleneck}
\label{subsec:tokenization}

Before any text reaches the Transformer layers, it is segmented into
subword units by a tokenizer, typically trained with byte-pair encoding
\citep{sennrich2016bpe} or a unigram language model as implemented in
SentencePiece \citep{kudo2018sentencepiece}. The tokenizer's vocabulary is
learned on (a sample of) the pre-training corpus, and therefore inherits
its language distribution: frequent English words receive dedicated tokens,
while words in under-represented languages are shattered into many short
fragments, in the worst case into individual bytes.

The severity of this effect is commonly quantified by \emph{fertility},
the average number of subword tokens produced per word (or per character
for scripts without whitespace word boundaries). High fertility has three
practical consequences. First, it inflates sequence lengths, making
inference disproportionately expensive for low-resource languages. Second,
it dilutes the semantic content carried by each token: a model that sees
Amharic or Tibetan only as byte sequences must reconstruct lexical meaning
from much noisier units than it does for English. Third, and most relevant
to this thesis, it affects the quality of sentence embeddings built from
token hidden states, since averaging over many weakly informative fragments
yields blurrier sentence representations. Non-Latin scripts (Ge'ez,
Myanmar, Khmer, Cherokee) and low-resource languages regardless of script
are the most affected. Throughout our experiments we therefore track
tokenizer fertility on both evaluation corpora alongside the alignment
scores, which allows us to distinguish genuinely weak cross-lingual
representations from representations that are handicapped already at the
tokenization stage.

% ..............................................
\subsection{Cross-lingual Transfer}
\label{subsec:transfer}

The practical value of multilingual models rests on \emph{cross-lingual
transfer}: capabilities acquired predominantly in one language (in
practice, English) become available in other languages without
language-specific supervision. The standard explanation is that the model
learns a partially \emph{shared semantic space} in which translation
equivalents --- words or sentences with the same meaning --- receive
similar representations, so that task knowledge attached to a region of
that space is accessible from any language that maps into it
\citep{hammerl2024survey}.

Empirical evidence supports this picture with an important refinement:
the degree of sharing is \emph{layer-dependent}. Input embeddings and the
first layers remain largely language-specific, dominated by lexical and
orthographic features; representations become increasingly
language-neutral through the middle of the network; and the final layers
re-specialize toward the output vocabulary of the target language.
Interpretability studies on Llama-family models show precisely this
trajectory \citep{wendler2024llamas,zhao2024multilingualism}, and it
motivates the layer-wise design of MEXA: if cross-lingual alignment exists
anywhere in an English-centric model, it should be found in the
intermediate layers, and a single-layer (e.g., last-layer) probe would
systematically underestimate it.

\citet{hammerl2024survey} organize the methods that \emph{create} or
\emph{improve} such alignment into families --- joint multilingual
pre-training, alignment objectives on parallel data, contrastive
sentence-level training, and post-hoc mapping between monolingual spaces.
This thesis is concerned with the complementary problem: not producing
alignment, but \emph{measuring} the alignment that a given model already
has, layer by layer and language by language.

% ------------------------------------------------------------
\section{English as a Pivot Language}
\label{sec:english-pivot}

% ..............................................
\subsection{Pivot Translation and Alignment}
\label{subsec:pivot-translation}

Using English as a bridge between languages long predates neural LLMs. In
statistical machine translation, direct parallel data between two
non-English languages is often scarce, so translation is routed through a
\emph{pivot} language: source $\rightarrow$ English $\rightarrow$ target
\citep{wu2007pivot}. Because English participates in the largest number of
parallel corpora, it is the pivot of choice, and the quality of any
source--target system becomes bounded by the quality of the two
English-anchored legs. Neural machine translation reproduced the pattern
in a softer form: multilingual NMT systems trained on English-centric data
develop an ``interlingua-like'' internal representation that enables
zero-shot translation between language pairs never seen together in
training \citep{johnson2017zeroshot}.

Modern decoder-only LLMs internalize this bridge. Using the logit-lens
technique, \citet{wendler2024llamas} show that when Llama-2 models process
non-English input, the intermediate layers pass through a phase in which
the representations are closest to \emph{English} tokens before the final
layers rotate them toward the target-language vocabulary --- the model
demonstrably ``works in English'' in its middle layers.
\citet{zhao2024multilingualism} corroborate this for reasoning tasks,
finding that multilingual inputs are converted into English-centric latent
representations in which the actual computation takes place. This
\emph{pivot hypothesis} is the theoretical foundation of MEXA: if
non-English understanding is routed through English-like intermediate
representations, then the degree to which a language's representations
align with English is a direct measurement of how much of the model's
(English-acquired) capability that language can access.

% ..............................................
\subsection{English Centricity and Bias}
\label{subsec:english-bias}

The pivot mechanism is a double-edged sword. On one side, it is exactly
what makes cross-lingual transfer from a 90\%-English corpus possible at
all. On the other, it builds a structural asymmetry into the model.
Languages that map cleanly onto English-like representations inherit the
model's full capability; languages that do not --- because of typological
distance, script, or sheer data scarcity --- are effectively locked out,
regardless of how much task knowledge the model possesses. Concepts and
distinctions that English does not encode (grammatical evidentiality,
honorific systems, culture-specific terminology) risk being flattened when
they are forced through an English-shaped bottleneck, and cultural framing
in generated text tends to default to Anglophone norms.

For evaluation, English centricity raises a methodological question that
this thesis addresses empirically: is English \emph{privileged} as a
measurement anchor, or merely \emph{convenient}? MEXA, as proposed, always
measures alignment against English. If the pivot hypothesis is the true
mechanism, replacing English with another pivot language should
systematically lower the measured alignment, in proportion to how far the
alternative pivot is from the model's internal hub. Our non-English-pivot
experiments (Section~\ref{sec:experimental-settings}) test exactly this by
recomputing alignment against French, German, Arabic, Chinese, and Basque
pivots.

% ------------------------------------------------------------
\section{Multilingual Evaluation}
\label{sec:multilingual-eval}

% ..............................................
\subsection{Standard Benchmarks}
\label{subsec:benchmarks}

The most direct way to assess multilingual capability is to run the model
on parallel downstream tasks. FLORES-200 \citep{nllbteam2022} provides
professionally translated Wikimedia sentences in 204 language--script
pairs and is the de-facto standard for evaluating translation and
cross-lingual retrieval. Belebele \citep{bandarkar2024belebele} builds a
machine reading comprehension benchmark on top of FLORES passages,
covering 122 language variants with fully parallel multiple-choice
questions --- crucially, differences in scores across languages can only
stem from the language itself, not from differing content. Knowledge- and
reasoning-oriented English benchmarks such as MMLU
\citep{hendrycks2021mmlu} and ARC have been machine-translated into
26--31 languages to form m-MMLU and m-ARC \citep{lai2023okapi}.

These benchmarks are indispensable, but they have structural limits: they
exist for at most a couple of hundred languages (Belebele: 122, m-MMLU:
34, m-ARC: 31), they are expensive to extend because each new language
needs high-quality translated task data, and machine-translated variants
inherit translation noise. For the several thousand remaining languages,
task-based evaluation is simply unavailable --- which is the gap that
alignment-based estimation aims to fill. In this thesis, downstream
benchmarks play the role of \emph{validation targets}: MEXA scores are
correlated against Belebele and m-ARC performance to test whether
alignment is a faithful proxy for task capability.

% ..............................................
\subsection{Reference-Based vs.\ Reference-Free Metrics}
\label{subsec:metrics}

Evaluation metrics in multilingual NLP fall into two broad families.
\emph{Reference-based} metrics compare model output against gold
references. Surface-overlap metrics such as BLEU
\citep{papineni2002bleu} are cheap and language-agnostic in principle, but
they reward n-gram overlap rather than meaning, penalize legitimate
paraphrases, and interact badly with morphologically rich languages and
non-whitespace scripts. Learned metrics such as COMET \citep{rei2020comet}
correlate far better with human judgments by scoring
source--hypothesis--reference triples with a fine-tuned multilingual
encoder, but they require human references and a trained scoring model,
and their reliability degrades precisely for the low-resource languages
where evaluation is most needed.

\emph{Reference-free} (quality-estimation or model-based) metrics remove
the dependency on gold outputs. Embedding-similarity approaches compare
representations of source and translation directly
\citep{reimers2020multilingual}, and model-internal approaches --- the
category MEXA belongs to --- do not look at generated text at all.
Instead, they probe the model's hidden states on parallel \emph{inputs}
and quantify structure that is known to correlate with capability. The
advantages are substantial: no task data, no references, no additional
trained judge, and applicability to any language for which a hundred
parallel sentences exist. The corresponding risk is validity --- an
internal signal is only useful insofar as it predicts external behavior
--- which is why correlation with downstream benchmarks
(Section~\ref{subsec:benchmarks}) is an essential part of the methodology
rather than an afterthought.

% ..............................................
\subsection{Challenges in Cross-lingual Alignment Evaluation}
\label{subsec:alignment-challenges}

Measuring alignment from raw embedding geometry is less straightforward
than it appears, for reasons rooted in the geometry of learned
representation spaces.

\paragraph{Anisotropy.} Transformer embedding spaces are strongly
anisotropic: representations occupy a narrow cone rather than spreading
over the full space, so the cosine similarity between two \emph{unrelated}
sentences is already high \citep{rajaee2022isotropy}. In multilingual
models the effect is compounded by outlier dimensions with extreme
variance that dominate cosine computations; removing a handful of such
dimensions can change similarity rankings entirely
\citep{hammerl2023anisotropy}. Absolute cosine values are therefore not
comparable across models, layers, or languages.

\paragraph{Hubness.} In high-dimensional spaces, some points become
``hubs'' that appear among the nearest neighbors of disproportionately
many queries. A naive nearest-neighbor criterion over cosine similarities
rewards such hubs and distorts retrieval-based measurements.

\paragraph{Language identity vs.\ semantics.} Embeddings encode both
\emph{what} a sentence says and \emph{which language} it is written in.
Two sentences in the same language are often closer to each other than a
sentence and its translation, simply because of the shared language
signal. A useful alignment metric must isolate the semantic component
from this language-identity component.

MEXA's design responds to all three problems at once: instead of using
cosine similarity as an absolute quantity, it uses it only
\emph{comparatively}, asking whether the correct translation outranks all
distractors within the same similarity matrix (a bidirectional
precision-at-1 criterion, formalized in
Section~\ref{sec:mexa-pipeline}). Rank-based comparisons within a fixed
matrix are invariant to the global scale and offset distortions that
anisotropy introduces, and the bidirectional requirement suppresses hub
sentences that would pass a one-directional test. As an extension, this
thesis additionally evaluates a \emph{mean-centered} variant in which each
language's mean embedding is subtracted before computing similarities,
directly removing the language-identity component discussed above.

% ------------------------------------------------------------
\section{Parallel Corpora}
\label{sec:parallel-corpora}

% ..............................................
\subsection{Sources and Typology}
\label{subsec:corpora-sources}

Alignment-based evaluation requires \emph{parallel} data: the same
sentences translated across many languages, so that semantic content is
held constant while language varies. Parallel corpora differ along
several axes that matter for measurement quality: translation quality
(professional vs.\ crawled/mined, e.g.\ OPUS \citep{tiedemann2012opus}),
domain (news, web, religious text), degree of parallelism (bitexts vs.\
$n$-way parallel corpora), and language coverage.

Two corpora sit at opposite ends of the coverage--quality trade-off and
are both used in this thesis. \textbf{FLORES-200} \citep{nllbteam2022}
contains professionally translated, quality-assured sentences from
Wikimedia sources in 204 language--script pairs; it is $n$-way parallel,
modern-domain, and high quality, but its coverage stops at roughly 200
languages. The \textbf{Parallel Bible Corpus} \citep{mayer2014bible}
exploits the fact that the Bible is the most translated text in existence,
with verse-level alignment across more than 1{,}400 language--script
pairs; the derived \emph{super-parallel} subset (sPBC) restricts to verses
available in all covered languages. Its domain is narrow and archaic and
translation styles vary, but it is the only resource that reaches deep
into the low-resource tail, which is exactly where reference-free
evaluation is most valuable.

% ..............................................
\subsection{Resource Disparity}
\label{subsec:resource-disparity}

The distribution of parallel (and monolingual) data across the world's
languages is extremely skewed. A few dozen high-resource languages account
for the overwhelming majority of available text; thousands of languages
have effectively no digital footprint beyond a Bible translation. This
disparity propagates through the entire pipeline studied in this thesis:
low-resource languages receive less pre-training data (weaker
representations), fragmented tokenization
(Section~\ref{subsec:tokenization}), and no downstream benchmarks
(Section~\ref{subsec:benchmarks}) --- so their capability can be neither
learned as well nor measured as easily.

For evaluation design, resource disparity has two concrete consequences.
First, any claim about ``multilingual'' capability that is validated only
on high-resource languages is systematically optimistic; extending
measurement to the sPBC's low-resource tail is therefore a substantive
robustness test, not an ornament. Second, corpus effects must be
disentangled from language effects: a language may score lower on the
Bible corpus than on FLORES because the model is genuinely weaker in it,
or because archaic religious register is far from the model's training
distribution. Using both corpora on the same models, as done here, makes
this confound visible and quantifiable.


% ============================================================
% CHAPTER 3 — METHODOLOGY
% ============================================================
\chapter{Methodology}
\label{ch:methodology}

This chapter specifies the experimental  of the thesis. We first
formalize the MEXA pipeline as implemented in our codebase
(Section~\ref{sec:mexa-pipeline}), then describe the evaluated models
(Section~\ref{sec:models}), the parallel corpora and their subsets
(Section~\ref{sec:datasets}), and finally the experimental settings,
analysis methods, and compute environment
(Section~\ref{sec:experimental-settings}). The methodology follows
\citet{kargaran2025mexa} for the reproduction experiments and extends it
along four axes: new decoder families, mixture-of-experts architectures,
encoder and embedding models, and non-English pivot languages.

% ------------------------------------------------------------
\section{The MEXA Pipeline}
\label{sec:mexa-pipeline}

MEXA (Multilingual Evaluation via Cross-Lingual Alignment) estimates the
multilingual coverage of a language model $m$ from parallel sentences
alone, in four stages: sentence-embedding extraction, layer-wise
similarity computation, alignment scoring, and aggregation.

\paragraph{Stage 1: Layer-wise sentence embeddings.}
Let $\mathcal{S} = \{s_1, \dots, s_n\}$ be a set of $n$ parallel sentences
available in a target language $L_1$ and a pivot language $L_2$ (English,
unless stated otherwise). Each sentence is passed through the model, and
hidden states $h_t^{(l)}$ are collected for every token $t$ at every layer
$l \in \{1, \dots, L\}$, including the embedding layer. Token states are
reduced to a single sentence embedding $e^{(l)}$ per layer with one of two
strategies:

\begin{itemize}
    \item \textbf{Position-weighted average.} Because causal attention
    means early tokens see little context while late tokens see all of it,
    tokens are averaged with weights proportional to their position:
    \begin{equation}
        e^{(l)} \;=\; \sum_{t=1}^{T} w_t\, h_t^{(l)},
        \qquad
        w_t \;=\; \frac{t}{\sum_{k=1}^{T} k},
        \label{eq:weighted-avg}
    \end{equation}
    where $T$ is the number of non-padding tokens. This is the primary
    setting in all reported results.
    \item \textbf{Last token.} The hidden state of the final non-padding
    token, $e^{(l)} = h_T^{(l)}$, which under causal attention is the only
    position that has attended to the entire sentence. This setting is
    used as a robustness check.
\end{itemize}

For encoder-only models (Section~\ref{subsec:encoder-models}), whose
attention is bidirectional, the same two reductions are applied to the
encoder hidden states, making the pipeline architecture-agnostic.
Embeddings are stored per language and per layer, so that scoring is
decoupled from (and far cheaper than) extraction.

\paragraph{Stage 2: Similarity matrix.}
For each layer $l$, a cosine similarity matrix
$C(L_1, L_2, m, l) \in \mathbb{R}^{n \times n}$ is computed between the
$n$ target-language embeddings and the $n$ pivot-language embeddings:
$c_{ij}(l)$ is the cosine similarity between target sentence $i$ and pivot
sentence $j$. The diagonal entries $c_{ii}(l)$ correspond to true parallel
pairs; all off-diagonal entries are non-parallel distractors.

\paragraph{Stage 3: Alignment score.}
As discussed in Section~\ref{subsec:alignment-challenges}, absolute cosine
values are unreliable in anisotropic LLM embedding spaces. MEXA therefore
scores alignment with a \emph{bidirectional retrieval} (precision-at-1)
criterion: a sentence pair counts as aligned only if its true similarity
strictly exceeds every distractor in both its row and its column,

\begin{equation}
    \mu C(L_1, L_2, m, l)
    \;=\;
    \frac{1}{n}
    \sum_{i=1}^{n}
    \mathbb{1}\!\left(
        c_{ii}(l) \;>\; \max_{j \neq i}\,\{\, c_{ij}(l),\; c_{ji}(l) \,\}
    \right).
    \label{eq:mexa}
\end{equation}

The score lies in $[0, 1]$; the expected value under random embeddings is
approximately $1/n^2$ per pair (both the row and the column condition must
hold by chance), so with $n = 100$ sentences a random baseline scores
essentially zero, which validates the use of small parallel samples.

\paragraph{Stage 4: Aggregation across layers.}
The layer-wise scores are summarized into two model-level statistics:
\begin{equation}
    \mu_{\text{Max}}(L_1) = \max_{l} \; \mu C(L_1, L_2, m, l),
    \qquad
    \mu_{\text{Mean}}(L_1) = \frac{1}{L} \sum_{l=1}^{L} \mu C(L_1, L_2, m, l).
\end{equation}
$\mu_{\text{Max}}$ captures the peak alignment a model achieves anywhere
in its depth --- typically in the middle-to-late layers, consistent with
the pivot hypothesis --- while $\mu_{\text{Mean}}$ rewards models that
sustain alignment across many layers. Model-level scores reported in
Chapter~4 average these quantities over all evaluated languages.

\paragraph{Extension: mean-centered scoring.}
To separate semantic alignment from the language-identity signal
(Section~\ref{subsec:alignment-challenges}), we additionally compute a
\emph{centered} variant of Equation~\eqref{eq:mexa}: before building the
similarity matrix at layer $l$, each language's embeddings are shifted by
their own mean,
\begin{equation}
    \tilde{e}_i^{(l)} \;=\; e_i^{(l)} - \frac{1}{n}\sum_{k=1}^{n} e_k^{(l)},
\end{equation}
applied independently to the target and pivot sides. Centering removes the
dominant direction shared by all sentences of a language and typically
sharpens retrieval; comparing centered and uncentered scores quantifies
how much measured alignment is carried by language-mean offsets rather
than sentence-level semantics.

\paragraph{Extension: non-English pivots.}
Equation~\eqref{eq:mexa} does not privilege English mathematically; the
pivot $L_2$ is a free parameter. We exploit this to test the pivot
hypothesis directly, recomputing all alignment scores with French, German,
Arabic, Chinese, and Basque as pivot languages --- spanning high-resource
Latin-script, non-Latin-script, and a language isolate --- in both the
standard and centered variants.

% ------------------------------------------------------------
\section{Models}
\label{sec:models}

We evaluate 25+ publicly available checkpoints, grouped into four
categories that correspond to the four experimental axes of the thesis.
All models are obtained from the Hugging Face Hub and evaluated without
any fine-tuning; base (non-instruction-tuned) variants are used where
available so that scores reflect pre-training rather than post-training.

% ..............................................
\subsection{Reproduction Models}
\label{subsec:reproduction-models}

To validate our implementation against \citet{kargaran2025mexa}, we
re-evaluate two models with published reference scores:
\textbf{Llama~3.1~8B} \citep{dubey2024llama3}, the primary English-centric
decoder studied in the original paper, and \textbf{Mistral~7B~v0.3}
\citep{jiang2023mistral}. Agreement with the published
$\mu_{\text{Max}}$/$\mu_{\text{Mean}}$ values on the Table-1
configurations (Section~\ref{sec:datasets}) establishes that our
embedding-extraction and scoring code reproduces the original pipeline,
and anchors all subsequent extensions.

% ..............................................
\subsection{New Decoder Families}
\label{subsec:new-decoders}

The first extension applies MEXA to decoder families released after the
original study, covering different data mixtures and scales:

\begin{itemize}
    \item \textbf{Qwen3} \citep{yang2025qwen3} at 0.6B, 1.7B, 4B, and 8B
    parameters --- a family trained with a substantially more multilingual
    data mixture than Llama, whose within-family size sweep isolates the
    effect of scale under a fixed recipe;
    \item \textbf{Qwen3.5~9B}, the successor generation, to measure
    recipe-level progress at approximately constant parameter count;
    \item \textbf{Apertus~8B} and \textbf{Apertus-Mini~4B}, openly trained
    models with an explicitly multilingual, compliance-focused data
    policy, providing a counterpoint to English-centric training;
    \item \textbf{Gemma~4} in five variants (E2B, E4B, 12B, 31B, and the
    sparse 26B-A4B), Google's current open-weight family, spanning
    efficient edge-scale to 31B-parameter models.
\end{itemize}

\noindent
Together with the reproduction models, this yields size ladders within
families (Qwen3 0.6B$\to$8B, Gemma~4 E2B$\to$31B) and recipe contrasts
across families at matched scale (Llama~3.1~8B vs.\ Qwen3~8B vs.\
Apertus~8B).

% ..............................................
\subsection{Mixture-of-Experts Models}
\label{subsec:moe-models}

Sparse mixture-of-experts (MoE) architectures route each token through a
small subset of expert feed-forward networks, decoupling parameter count
from per-token compute \citep{jiang2024mixtral}. Whether expert routing
helps or hurts cross-lingual alignment is an open question: experts could
specialize by language (fragmenting the shared space) or by function
(leaving alignment intact). We evaluate \textbf{Mixtral~8x7B} and
\textbf{Mixtral~8x22B} \citep{jiang2024mixtral}, comparing them against
their dense sibling Mistral~7B~v0.3, and \textbf{Gemma~4~26B-A4B},
comparing it against dense Gemma~4 variants. In all cases the MEXA
pipeline applies unchanged, since it only consumes the post-layer hidden
states.

% ..............................................
\subsection{Encoder and Embedding Models}
\label{subsec:encoder-models}

The final model axis contrasts causal decoders with bidirectional
encoders and purpose-built sentence-embedding models, which provide an
upper-bound reference for what alignment looks like when it is an explicit
training objective rather than a by-product:

\begin{itemize}
    \item \textbf{Masked-LM encoders:} XLM-R base and large
    \citep{conneau2020xlmr}, Glot500 \citep{imani2023glot500} (an XLM-R
    variant extended to 500+ languages, particularly relevant for the
    Bible corpus tail), and mmBERT-base;
    \item \textbf{Contrastively trained sentence encoders:} LaBSE
    \citep{feng2022labse} and multilingual-E5-base
    \citep{wang2024multilingual}, both explicitly optimized so that
    translations are nearest neighbors;
    \item \textbf{Decoder-based embedding models:} Qwen3-Embedding at
    0.6B, 4B, and 8B, which share the Qwen3 backbone but are fine-tuned
    for retrieval --- allowing a controlled comparison of the same
    architecture before and after embedding-specific training.
\end{itemize}

% ------------------------------------------------------------
\section{Datasets}
\label{sec:datasets}

All experiments draw on the two parallel corpora introduced in
Section~\ref{sec:parallel-corpora}, in configurations chosen to match the
original paper where reproduction is intended and to extend coverage
elsewhere.

\paragraph{FLORES-200.}
From the FLORES-200 devtest split \citep{nllbteam2022} we use:
\begin{itemize}
    \item \textbf{FLORES Table-1 subset:} the 116 languages overlapping
    with Belebele, with $n = 100$ parallel sentences per language ---
    the exact configuration of Table~1 in \citet{kargaran2025mexa}, used
    for reproduction and for all model-family comparisons;
    \item \textbf{FLORES Table-1 (2000):} the same 116 languages with all
    $n = 2{,}000$ devtest sentences, to quantify the effect of sample
    size on the retrieval criterion (larger $n$ means more distractors
    and therefore a strictly harder task);
    \item \textbf{FLORES Full:} all 204 language--script pairs, for
    maximal high-quality coverage.
\end{itemize}

\paragraph{Parallel Bible Corpus (sPBC).}
From the super-parallel Bible corpus \citep{mayer2014bible} we use:
\begin{itemize}
    \item \textbf{Bible Table-1 subset:} 101 languages with $n = 103$
    parallel verses, matching the original paper's Bible configuration;
    \item \textbf{Bible Full:} the complete corpus of more than 1{,}400
    language--script pairs, which extends evaluation far into the
    low-resource tail where no downstream benchmark exists.
\end{itemize}

English (\texttt{eng\_Latn}) serves as the pivot in all standard runs; the
pivot-substitution runs of Section~\ref{sec:mexa-pipeline} use the FLORES
Table-1 subset. Language identifiers follow the FLORES-200 convention of
ISO~639-3 code plus script (e.g.\ \texttt{shn\_Mymr}).

% ------------------------------------------------------------
\section{Experimental Settings}
\label{sec:experimental-settings}

% ..............................................
\subsection{Evaluation and Analysis Methods}
\label{subsec:analysis-methods}

Unless stated otherwise, reported scores use position-weighted average
embeddings (Equation~\eqref{eq:weighted-avg}) and are given as
$\mu_{\text{Max}}$ and $\mu_{\text{Mean}}$ per model--dataset
configuration; last-token embeddings and the centered variant are reported
as ablations. The analysis proceeds on four levels:

\begin{itemize}
    \item \textbf{Model level:} scores averaged over all languages of a
    configuration, enabling family- and scale-comparisons and the
    reproduction check against published values;
    \item \textbf{Language level:} per-language scores, ranked and grouped
    by script and by resource level, and cross-referenced with tokenizer
    fertility (computed on both corpora) to separate tokenization effects
    from representational effects;
    \item \textbf{Layer level:} full layer-wise alignment profiles
    $\mu C(\cdot, l)$, visualized as heatmaps, to locate where in the
    network alignment peaks and to test the intermediate-layer prediction
    of the pivot hypothesis;
    \item \textbf{Validity level:} Pearson correlations between MEXA
    scores and downstream accuracy on Belebele
    \citep{bandarkar2024belebele} and m-ARC \citep{lai2023okapi} over the
    languages both cover, testing whether alignment is a faithful proxy
    for task performance; cross-model Pearson correlations of per-language
    score profiles additionally measure how consistent alignment
    structure is across architectures.
\end{itemize}

For qualitative inspection, sentence embeddings are also projected to two
and three dimensions with PCA and t-SNE, colored by language and script,
to visualize the clustering behavior that the retrieval score summarizes
numerically.

% ..............................................
\subsection{Implementation and Compute}
\label{subsec:implementation}

The pipeline is implemented in Python with PyTorch and Hugging Face
Transformers. \texttt{embed\_extractor.py} loads each checkpoint via
\texttt{AutoModelForCausalLM} (or \texttt{AutoModel} for encoders) with
output of all hidden states enabled, tokenizes each sentence with the
model's own tokenizer, and stores per-layer sentence embeddings as
serialized NumPy arrays, one file per language.
\texttt{compute\_mexa.py} then builds the layer-wise cosine similarity
matrices in float64 and evaluates Equation~\eqref{eq:mexa}; because
scoring operates on cached embeddings, ablations over pooling strategies,
pivots, and centering re-use a single (expensive) extraction pass per
model--corpus pair. Formatted per-language and per-layer scores are
exported to CSV/JSON and explored in an interactive React dashboard
developed for this thesis.

Small models were run locally; all large-scale extraction jobs (8B and
above, MoE, and full-corpus sweeps) were executed on the CIT-TUM-HN
SLURM cluster at TUM Campus Heilbronn on NVIDIA data-center GPUs, with
models loaded in reduced precision (bfloat16, with quantized loading for
the largest MoE checkpoints) and embeddings cast to float32 before
storage. Random effects play no role in the pipeline --- extraction and
scoring are deterministic given a checkpoint and corpus --- so no seed
averaging is required.
